\documentclass[11pt]{amsart}
\usepackage[T1]{fontenc}
\usepackage{lmodern}
\usepackage{amsmath,amssymb,mathtools,booktabs,array}
\usepackage{microtype}
\usepackage[hidelinks]{hyperref}
\usepackage[margin=1.05in]{geometry}
\newtheorem{theorem}{Theorem}[section]
\newtheorem{lemma}[theorem]{Lemma}
\newtheorem{proposition}[theorem]{Proposition}
\newtheorem{corollary}[theorem]{Corollary}
\theoremstyle{definition}
\newtheorem{definition}[theorem]{Definition}
\theoremstyle{remark}
\newtheorem{remark}[theorem]{Remark}
\newcommand{\N}{\mathbb N_0}
\newcommand{\Fib}{\mathsf F}
\DeclareMathOperator{\supp}{supp}
\numberwithin{equation}{section}
\title{Genus monotonicity for numerical semigroups}
\author{Mingchang Liu}
\email{mliu416@gatech.edu}
\subjclass[2020]{Primary 20M14; Secondary 05A15, 05A20}
\keywords{Numerical semigroup, genus, Kunz coordinates, generating function}
\date{}
\begin{document}
\begin{abstract}
Let $n_g$ denote the number of numerical semigroups of genus $g$.
We prove that $n_{g+1}\ge n_g$ for every $g\ge0$, settling
the weak monotonicity form of the Bras-Amor\'os conjecture.
The computer-assisted proof constructs an injection outside a specified family of leaves
and bounds this exceptional family by a collection of unused targets.
Generating-function estimates establish the comparison for $g\ge837$;
exact transfer computations bound relaxed counts of the exceptional
sources for $1\le g\le836$.
The analytic estimates distinguish genuine semigroups from a larger
two-pivot relaxation, whose coefficients can be computed by transfer
along paths and cycles.
\end{abstract}
\maketitle

\section{Introduction}

A numerical semigroup is an additive submonoid $S$ of $\N$ with finite
complement. Its genus is $g(S)=|\N\setminus S|$. Write
\[
 \mathcal S_g=\{S:g(S)=g\},\qquad n_g=|\mathcal S_g|.
\]
Bras-Amor\'os's computations led her to conjecture the stronger Fibonacci
inequality for this sequence~\cite{Bras}. We consider its weak monotonicity
form, recorded separately in Kaplan's survey~\cite[Conjecture~2]{Kaplan}.

\begin{theorem}\label{thm:main}
For every integer $g\ge0$, one has $n_{g+1}\ge n_g$.
\end{theorem}

The conjecture concerns all numerical semigroups, with no restriction on
multiplicity or conductor. It is weaker than the Fibonacci inequality
$n_g\ge n_{g-1}+n_{g-2}$ for $g\ge2$, which is not established here.
The distinction between all genera and sufficiently large genus is essential.
Zhai proved that
\begin{equation}\label{eq:known-asymptotic}
 n_g\sim C\varphi^g,\qquad
 \varphi=\frac{1+\sqrt5}{2},\quad C>0.
\end{equation}
Consequently $n_{g+1}/n_g\to\varphi>1$, so eventual strict increase was
already known~\cite{Zhai}. An all-genus proof requires control of the
range preceding that asymptotic conclusion. Exact values of $n_g$
have been computed through genus $80$~\cite[Section~4]{BT}.

The argument starts with an incomplete injection rather than an estimate
of $n_g$ itself. For each $g$, we construct an injection from
$\mathcal S_g\setminus\mathcal E_g$ into $\mathcal S_{g+1}$ and exhibit
at least $L(g)$ targets outside its image. Thus
\begin{equation}\label{eq:master-inequality}
 n_{g+1}-n_g\ge L(g)-|\mathcal E_g|.
\end{equation}
The exceptional sources are leaves in the numerical-semigroup tree and
have depth four or five. The construction is arranged so that their
remaining possibilities admit two different estimates: a coefficient
bound for a finite range and a weighted generating-function bound for
large genus. Neither estimate is required to approximate $n_g$ closely.

The proof uses established low-depth constructions. Eliahou and Fromentin
constructed two disjoint injections on gapsets of depth at most three,
which become appending $1$ and appending $2$ in Kunz coordinates~\cite[Propositions~5.2--5.3]{EF}.
Their results imply a Fibonacci lower recurrence for the low-depth count.
Kunz coordinates and their interpretation as words or compositions also
underlie the asymptotic work of Bacher and Zhu~\cite{Bacher,Zhu}.
Bacher uses maximal pivots, reflections and transfer matrices; Zhu proves
sub-Fibonacci exponential bounds for the higher-depth population using
weighted graph-homomorphism estimates. The general use of these methods
is therefore established. Here they are combined with a specified
exceptional family, disjoint unused targets, and explicit constants that
permit the finite and analytic ranges to meet.

The all-genus monotonicity problem is distinguished explicitly from
its stronger variants in recent work of Cyrusian and Kaplan~\cite{CK}.
Bras-Amor\'os and Torra obtain refined recurrences under multiplicity
and Frobenius-number hypotheses~\cite{BT}; Casas, Madrid and Rosales
study related internal and leaf counts~\cite{CMR}.
For semigroups of ordinarization number three, Kumar's preprint combines
a shift on a large class with a rank matching into a disjoint reserve,
using separate finite and eventual capacity bounds~\cite[Sections~7--8]{Kumar}.
This gives a restricted precedent for the allocation strategy used here.
An earlier proof of monotonicity was proposed by Kachouei and
Rahmati~\cite{KR}. The proof of their Theorem~19 (p.~13) asserts
\[
 \binom{n_{m,g-1}}{m-1}\le n_{m,g}\qquad(n_{m,g-1}\ge m),
\]
where $n_{m,g}$ counts semigroups of multiplicity $m$ and genus $g$.
However, $n_{3,9}=n_{3,10}=4$, so the bound would require
$6=\binom42\le4$.
Indeed, multiplicity-three Kunz pairs satisfy
$a+b=g$, $2a\ge b$ and $2b+1\ge a$, giving these counts directly.
This concerns the proposed counting argument, not the monotonicity
statement itself.

Our finite computation evaluates a relaxed counting problem and checks
the inequalities in~\eqref{eq:master-inequality}. It does not enumerate
all numerical semigroups through genus $836$, or compute all the
corresponding values of $n_g$. We give the coverage and arithmetic
justification separately from the algorithms' implementation.
The proof of Theorem~\ref{thm:main} is completed after the finite and
analytic comparisons have been established.


The population notation used throughout the proof is collected in
Table~\ref{tab:notation}. Auxiliary matrix entries and scalar bounds are
local to the subsection in which they are defined.

\begin{table}[ht]
\centering
\caption{The principal populations and counting bounds.}
\label{tab:notation}
\begin{tabular}{@{}lp{.72\linewidth}@{}}
\toprule
Notation & Meaning\\
\midrule
$\mathcal S_g$, $n_g$ & All genus-$g$ semigroups and their number\\
$\mathcal E_g$ & Sources left unmatched by the partial injection\\
$L(g)$ & Explicit lower bound on unused genus-$(g+1)$ targets\\
$A_g$, $B_g$ & Genuine words with at least two maxima, of depth four and five\\
$A(z)$, $B(z)$ & Generating series of these two genuine populations\\
$Z_{m,r,p}(z)$ & Raw depth-four partition with rightmost fours at $r<p$\\
$F_{r,p}(z)$ & Sum of $Z_{m,r,p}(z)$ over $m>p$\\
$U_g$, $V_g$ & Integer upper bounds used in the finite comparison\\
\bottomrule
\end{tabular}
\end{table}



\section{A partial injection and its exceptional sources}
\label{cmb:injection}

We use positive Kunz coordinates~\cite{Kunz,Rosales}. If a numerical semigroup $S$ has positive
genus and multiplicity $m$, let $i+mx_i$ be its least element congruent to
$i$ modulo $m$, for $1\le i<m$. Its Kunz word is
$x=(x_1,\ldots,x_{m-1})$. The empty word represents $\mathbb Z_{\ge0}$.
The weight of a word is $|x|=\sum_i x_i$; its depth is $\max_i x_i$
for a nonempty word. The empty word has weight and depth zero.
A word is called binary if all its letters belong to $\{1,2\}$.
Throughout, coefficients with negative indices are zero.

\begin{lemma}\label{cmb:kunz}
Positive words of length $m-1$ represent numerical semigroups of
multiplicity $m$ precisely when
\begin{align}
 x_i+x_j&\ge x_{i+j} &&(i+j<m),\label{cmb:noncarry}\\
 x_i+x_j+1&\ge x_{i+j-m} &&(i+j>m),\label{cmb:carry}
\end{align}
including self-pairs. The genus is $|x|$. If the depth is $a$ and $p$ is the
rightmost coordinate equal to $a$, the Frobenius number is $(a-1)m+p$.
\end{lemma}
\begin{proof}
In residue $i$, the gaps are $i,i+m,\ldots,i+(x_i-1)m$, proving the genus
formula. Adding the least nongaps in two nonzero residues proves necessity
of the displayed inequalities. Conversely, these inequalities ensure that
sums of these least nongaps belong to the prescribed residue tails. Adding
multiples of $m$ then proves closure for all nongaps. A sum in residue zero
already belongs to the semigroup, so no condition is needed when $i+j=m$.
The largest gap in residue $i$ is $i+(x_i-1)m$, and comparison of these
numbers proves the Frobenius formula.
\end{proof}

We set the Frobenius number of $\mathbb Z_{\ge0}$ to $-1$.
An \emph{effective generator} is a minimal generator exceeding the Frobenius
number. Deleting one gives an ordinary child of genus one greater. A
semigroup without effective generators is a \emph{leaf}. We call depth at
most three low depth and depth at least four high depth.

\begin{lemma}\label{cmb:parent}
Every positive-genus semigroup has a unique ordinary parent, obtained by
restoring its Frobenius number. Every parent has at most one leaf child.
For a high-depth word, restoring the Frobenius decreases its rightmost
maximal coordinate by one and preserves multiplicity. Consequently a
high-depth leaf has a low-depth parent exactly when its depth is four and
it has a unique coordinate equal to four.
\end{lemma}
\begin{proof}
Let $F$ be the Frobenius number. Restoring $F$ preserves closure: every sum
of $F$ and a positive nongap is larger than $F$, as is $2F$. In the enlarged
semigroup, $F$ cannot be a sum of two positive elements, for such a
decomposition would already have contradicted closure before restoration.
The new Frobenius is smaller than $F$, so $F$ is an effective generator.
Conversely, deletion of an effective generator $a$ gives a semigroup whose
Frobenius is $a$, recovering its parent uniquely.

If $a<b$ are effective generators of a parent, deleting $a$ leaves $b$
minimal, because deletion cannot introduce a decomposition. It also leaves
$b$ effective, because the child's Frobenius is $a$. Thus only deletion of
the largest effective generator can give a leaf. At high depth, $F>m$;
restoration and deletion of effective generators therefore preserve $m$.
The coordinate statement follows from Lemma~\ref{cmb:kunz}, and yields the
last assertion. At multiplicity two a word $(k)$ has effective generator
$2k+1$; in particular this boundary introduces no high-depth leaves.
\end{proof}

All sets at a fixed genus are finite. Indeed the map $s\mapsto F-s$
injects the nongaps in $[0,F]$ into its gaps, whence $F\le2g-1$ at positive
genus $g$. Ordering the finite gap sets lexicographically, for example by
their indicator vectors in $[1,2g-1]$, supplies an order whenever we use a
finite-set matching below.

Fix $g\ge1$. Partition the set $\mathcal S_g$ of genus-$g$ semigroups by
the following tests, applied in order; each test selects only sources not
selected earlier:
\begin{enumerate}
\item $\mathcal L_g$: all low-depth words;
\item $\mathcal P_g$: all high-depth nonleaves;
\item $\mathcal C_g$: all remaining high-depth leaves with a low-depth parent;
\item $\mathcal T_g$: all remaining words $(5,u)$ with every letter of $u$
  at most three;
\item $\mathcal D_g$: all remaining words of depth at least six, together
  with all remaining depth-four or depth-five words containing no $1$;
\item $\mathcal I_g$: all remaining words with a unique $5$, at a position
  at least two, and all other letters at most three;
\item $\mathcal B_g$: all remaining depth-four words whose sole $1$ is final;
\item $\mathcal E_g$: all remaining sources.
\end{enumerate}

\begin{lemma}\label{cmb:residual}
The exceptional set $\mathcal E_g$ is precisely the set of leaves containing
at least one $1$ in the following three disjoint sectors:
\begin{enumerate}
\item depth four with at least two $4$'s, excluding words whose sole $1$ is final;
\item depth five with a unique $5$ and at least one $4$;
\item depth five with at least two $5$'s.
\end{enumerate}
\end{lemma}
\begin{proof}
The first two tests leave only high-depth leaves. By Lemma~\ref{cmb:parent},
the third removes exactly the depth-four leaves with a unique $4$.
The fifth removes all depths at least six and all words without $1$.
A depth-five word with only one letter at least four has a unique $5$;
it belongs to the fourth test if this is its first letter, and otherwise
to the sixth. The seventh test makes the stated further exclusion at depth
four. Every leaf in one of the three displayed sectors misses all seven
tests, proving equality.
\end{proof}

The append-$1$ and append-$2$ maps on low-depth words are the
Eliahou--Fromentin injections~\cite[Propositions~5.2--5.3]{EF}. We recall their elementary proof
because the change of multiplicity matters here. Appending either letter
to a word of multiplicity $m$ makes the new multiplicity $m+1$. Old
noncarry inequalities remain unchanged. The new noncarry target, at index
$m$, is at most two; all its inequalities are automatic. Every new carry
left side is at least three, and every target is at most three. Thus both
append operations preserve validity.

Define the following maps:
\begin{equation}\label{cmb:maps}
\begin{array}{rcll}
 x&\longmapsto&(x,1),&x\in\mathcal L_g,\\
 x&\longmapsto&\text{the child obtained by deleting the largest effective generator},
      &x\in\mathcal P_g,\\
 x&\longmapsto&(\operatorname{par}(x),2),&x\in\mathcal C_g,\\
 (5,u)&\longmapsto&(2,u,1,3),&(5,u)\in\mathcal T_g.
\end{array}
\end{equation}
Here $\operatorname{par}(x)$ denotes the parent's Kunz word. For
$x=(w,1)\in\mathcal B_g$, replace each letter of $w$ using
\begin{equation}\label{cmb:code}
 2\longmapsto11,\qquad3\longmapsto21,\qquad4\longmapsto22,
\end{equation}
and append $2$.

\begin{lemma}\label{cmb:elementary-maps}
These maps are injective into $\mathcal S_{g+1}$ and their image sectors
are pairwise disjoint. The $\mathcal B_g$ image consists of binary words
ending in $2$.
\end{lemma}
\begin{proof}
The append-$1$ map is valid and reversible. A source in $\mathcal C_g$
has a depth-three parent of genus $g-1$; appending $2$ therefore gives a
nonbinary word of genus $g+1$. Removing this final letter determines the
parent, and Lemma~\ref{cmb:parent} determines its leaf child uniquely.
The ordinary-child map is valid, has high depth, and is inverted by
restoration of the Frobenius number.

For the fourth map, suppose the original multiplicity is $m$. The new
multiplicity is $m+2$, and the weight changes by $-5+2+1+3=1$.
Old noncarry targets have value at most three. If an addend has changed,
its new value is $2$, so its sum with any positive letter is sufficient.
The target at index $m$ is $1$. For the final target $3$ at index $m+1$,
a pair of old indices had, in the source, a carry inequality targeting the
initial $5$, so its old sum was at least four. If one of these indices were
$1$, the other would be $m$, which is not an old index. Thus these sums
are unchanged. The sole additional pair is $(1,m)$, of values $(2,1)$.
Every carry inequality is automatic at depth three. Removing the final
$1,3$ and changing the initial $2$ back to $5$ gives the inverse.

In the last map, $w$ is a word over $\{2,3,4\}$ with at least two $4$'s.
Every block preserves weight, and replacement of the final $1$ by $2$
raises genus by one. Binary words satisfy all Kunz inequalities. Remove
the final $2$ and parse consecutive pairs to recover $w$ uniquely.

The first four images are respectively low depth ending in $1$, high
depth, low depth and nonbinary ending in $2$, and low depth ending in $3$.
The code image is binary ending in $2$. These descriptions are disjoint.
\end{proof}

\section{Binary allocation}
\label{cmb:binary}
Write $\Fib_0=0$, $\Fib_1=1$, and
$\Fib_{g+1}=\Fib_g+\Fib_{g-1}$ for $g\ge1$.
There are exactly $\Fib_g$ binary words of weight $g+1$ ending in $2$:
remove the final $2$ and count compositions of $g-1$ over $\{1,2\}$.
This includes $g=1$, when the remaining composition is empty. We bound
jointly the sources assigned to this sector, including its code images.

Define the following nonnegative formal series:
\begin{align*}
 a_k(z)&=\sum_{\lceil k/2\rceil\le a\le k}z^a,&
 b_k(z)&=\sum_{\substack{1\le a,b\le k\\a+b\ge k}}z^{a+b},&
 K_k(z)&=\frac{1+a_k(z)}{1-b_k(z)},\\
 a(z)&=\frac{z^4}{1-z},&
 b(z)&=\frac{z^8(7-6z)}{(1-z)^2}.
\end{align*}
Put
\begin{align}
 M(z)&=z^6K_6K_5+z^7K_7K_6+z^8K_8K_7
       +\frac{z^9(1+a)^2}{(1-z)(1-b)^2},\label{cmb:M}\\
 H(z)&=\frac{z^4+z^5}{(1-z^2-z^3)(1-z^2-z^3-z^4-z^5)},\nonumber\\
 J(z)&=\frac{2z^{12}+4z^{14}+2z^{15}}
 {(1-2z^5-z^6)(1-z^2-z^3)(1-3z^4-2z^5-z^6)},\nonumber\\
 K(z)&=\frac{z^9}{(1-z^2-z^3)^2(1-z^2-z^3-z^4)},\nonumber\\
 W(z)&=M(z)+H(z)+J(z)+K(z),\qquad W_g=[z^g]W(z).\label{cmb:W}
\end{align}
The notation $K(z)$ and $K_k(z)$ denotes different functions.

\begin{lemma}\label{cmb:majorant}
For every $g\ge1$,
\[
 |\mathcal D_g|+|\mathcal I_g|+|\mathcal B_g|\le W_g.
\]
\end{lemma}
\begin{proof}
At the rightmost maximum $q$, the coordinates to its left are paired by
reflection about that position: each pair has sum at least $q$, and its
letters are at most $q$. To its right, each pair has sum at least $q-1$
and its letters are at most $q-1$. A central self-pair has the corresponding
ceiling-half lower bound. For either block, summing over even and odd
lengths gives $K_q$ or $K_{q-1}$. The two block lengths determine the pivot
position, so $z^qK_qK_{q-1}$ bounds all words of depth $q$ without an
additional position factor. For $q\ge9$, allow arbitrary positive letters,
weaken both pair thresholds to eight and both central thresholds to four.
The central series is $a$, and the ordered pair series is
$\sum_{s\ge8}(s-1)z^s=b$. Summing the pivot weight over $q\ge9$ yields
the last term of~\eqref{cmb:M}. Thus $M$ bounds depth at least six.
Subtracting the composition series on $\{2,3\}$ from that on
$\{2,3,4,5\}$ gives $H$, which counts all no-$1$ words of depth four or
five. These two estimates bound $\mathcal D_g$.

For an isolated $5$ at a noninitial position, the left block is nonempty,
uses letters at most three, and has reflection-pair sum at least five.
In particular it contains no $1$. Its series is $L_0-1$, where
\[
 L_0=\frac{1+z^3}{1-2z^5-z^6}.
\]
The right block has pair threshold four and central threshold two, giving
\[
 Q_0=\frac{1+z^2+z^3}{1-3z^4-2z^5-z^6}.
\]
Since the no-$1$ words were already removed, the right block must contain
$1$. Its subpopulation without $1$ is exactly all $\{2,3\}$ words, with
series $C_0=(1-z^2-z^3)^{-1}$. Direct simplification gives
$z^5(L_0-1)(Q_0-C_0)=J$, proving the required bound for $\mathcal I_g$.

Finally, splitting a $\{2,3,4\}$ word at its first two $4$'s gives the series
\begin{equation}\label{cmb:N}
 N(z)=\frac{z^8}{(1-z^2-z^3)^2(1-z^2-z^3-z^4)}.
\end{equation}
A final $1$ supplies the factor $z$, so $K=zN$ majorizes $\mathcal B_g$.
We discarded inequalities only to enlarge each population. Summing the
three bounds proves the assertion.
\end{proof}

\begin{lemma}[Binary numerical bound]\label{cmb:binary-numerical}
For all $g\ge0$, $W_g\le13\Fib_g/25$. For $g\ge141$,
$W_g\le13\Fib_g/1000$. Moreover, with $r=809/500$,
\begin{equation}\label{cmb:binary-decay}
 \frac{W_g}{\Fib_g}<\frac{50r^2}{((63/100)r)^g}\qquad(g\ge1).
\end{equation}
\end{lemma}
\begin{proof}
The exact finite arithmetic inputs are
\begin{align*}
 25W_g&\le13\Fib_g &&(0\le g\le499),\\
 1000W_g&\le13\Fib_g &&(141\le g\le499),\\
 W(63/100)&<50,&
 \frac{50r^2}{((63/100)r)^{500}}&<\frac{13}{1000}.
\end{align*}
All geometric denominators in~\eqref{cmb:W} are positive at $63/100$.
These are finite rational checks. In particular, for each rational summand
$P(z)/Q(z)$, with $Q(0)=1$, its coefficients are calculated using
\begin{equation}\label{cmb:coefficient-recurrence}
 f_n=[z^n]P-\sum_{j=1}^{\min(n,\deg Q)}[z^j]Q\,f_{n-j}.
\end{equation}
This recurrence and the displayed rational formulas specify the exact
arithmetic checks without numerical rounding.
Since $r^2<r+1$, the two initial cases and induction in the Fibonacci
recurrence give $\Fib_g\ge r^{g-2}$ for $g\ge1$. Nonnegativity gives
$W_g(63/100)^g\le W(63/100)<50$, proving~\eqref{cmb:binary-decay}.
The base $(63/100)r$ is greater than one, so the sharper finite bound
extends to every $g\ge500$ and also implies the weaker bound there.
\end{proof}

Reserve the actual images~\eqref{cmb:code} first. In the fixed finite-set
orders, match $\mathcal D_g\cup\mathcal I_g$ to the binary end-$2$ words
not reserved. Lemmas~\ref{cmb:majorant} and~\ref{cmb:binary-numerical}
show that the total demand, including reserved images, is at most
$13\Fib_g/25$, hence at most the supply $\Fib_g$.
This matching is injective and its images are disjoint from all images
in Lemma~\ref{cmb:elementary-maps}. We have proved the following.

\begin{proposition}\label{cmb:partial}
There is an injection
$\Phi_g:\mathcal S_g\setminus\mathcal E_g\longrightarrow\mathcal S_{g+1}$.
At least
\begin{equation}\label{cmb:binary-capacity}
 L_{\rm bin}(g)=\max\left\{\left\lceil\frac{12\Fib_g}{25}\right\rceil,
                                      \Fib_g-W_g\right\}
\end{equation}
binary targets of genus $g+1$ ending in $2$ lie outside its image.
\end{proposition}
The two entries in this maximum bound the same unused set. In particular,
they must not be added. For $g\ge141$, this gives
$L_{\rm bin}(g)\ge987\Fib_g/1000$.

\section{Witnessed unused targets}
\label{cmb:reservoir}
Let a prefix $P$ of length $n\ge1$, ending in $3$, be followed by a binary
tail of length $\ell$. The base word has multiplicity $m=n+\ell+1$.
A prefix $3$ at position $i=n-d$ is \emph{witnessed} if two tail positions
$n+j,n+k$, both containing $1$, satisfy
\begin{equation}\label{cmb:witness}
 j+k=\ell+1-d.
\end{equation}
We permit $j=k$. This is exactly the carry equality
$(n+j)+(n+k)=m+i$. Hence replacing that $3$ by $4$ violates its carry
inequality, whose left side is $1+1+1=3$.

\begin{lemma}\label{cmb:witness-unused}
If a base word is valid and every $3$ is witnessed, appending $2$ gives a
nonbinary target outside the image of $\Phi_g$, where $g$ is one greater
than the weight of the base. Distinct bases give distinct unused targets.
\end{lemma}
\begin{proof}
The base has depth three, so its Frobenius exceeds its multiplicity. Its
ordinary children preserve multiplicity, and a high-depth child would
have to increase a coordinate $3$ to $4$. Every such increase is blocked
by~\eqref{cmb:witness}, so no such child exists. Appending $2$ preserves
validity. A target in the $\mathcal C_g$ image would recover, on removing
this $2$, a base having a high-depth leaf child, a contradiction. Depth,
final letter, and the presence of $3$ exclude every other image sector.
Deleting the appended $2$ recovers the base. Its weight is $g-1$, while
the target weight is $g+1$, giving the stated genus shift.
\end{proof}

Put $u=z^2+z^3$, $s=z+z^2+z^3$, $c=z+z^2$, and $v=us$.
We use two disjoint families of prefixes. A \emph{safe-half prefix} of
length $n$ has its first $\lfloor n/2\rfloor$ letters in $\{2,3\}$,
its remaining nonfinal letters in $\{1,2,3\}$, and its last letter $3$.
This is Bacher's family in the proof of~\cite[Theorem~15.1, Section~15.1]{Bacher}.
We use it here with tail witnesses to obtain targets outside the partial
injection; the prefix family itself is not new.
An \emph{early-single-one prefix} has its unique prefix $1$ at a position
$p\ge1$, has $2$ at position $2p$, ends in $3$ at some $n\ge2p+1$,
and has all other letters in $\{2,3\}$.

Every base obtained from either family and an arbitrary binary tail is
valid. At depth three only a noncarry $1+1$ targeting $3$ can fail. In the
safe-half family two prefix $1$ positions sum beyond the prefix. In the
early-single-one family the only prefix pair of $1$'s is the self-pair at
$p$, targeting the prescribed $2$ at $2p$. A pair involving a tail position
cannot noncarry-target the prefix. These observations prove validity.
The early $1$ lies in the first half, so the two families are disjoint.
The last $3$ uniquely identifies the prefix in a base.

Writing $U=\{2,3\}$ and $V=\{1,2,3\}$, the safe prefixes have forms
$U^kV^k3$ and $U^{k+1}V^k3$, where powers denote independently chosen
letters. The early prefixes have form $U^h1U^h2U^k3$; the two $U^h$ blocks
need not be identical. Therefore their weight series are
\begin{equation}\label{cmb:prefix-series}
 C_s=\frac{z^3(1+u)}{1-v},\qquad
 C_e=\frac{z^6}{(1-u)(1-u^2)}.
\end{equation}

We now count bases with at least one unwitnessed $3$ by marking a failed
$3$. Let $t$ record the distance of the mark from the prefix end. The two
marked-prefix series are
\begin{align}
 F_s(z,t)&=C_s+
 \frac{z^6u(1+u)t}{(1-v)(1-vt)}+
 \frac{z^6t(1+st)}{(1-vt)(1-vt^2)},\label{cmb:marked-safe}\\
 F_e(z,t)&=C_e+
 \frac{z^9t}{(1-u^2)(1-u)(1-ut)}+
 \frac{z^9ut^2}{(1-ut)(1-u^2)(1-u^2t)}
 \nonumber\\
 &\quad+\frac{z^9ut^4}{(1-ut)(1-u^2t)(1-u^2t^2)}.
 \label{cmb:marked-early}
\end{align}
To verify these formulas directly, marking the terminal $3$ contributes
$C_s$ or $C_e$. In a block of length $k$ with letter-weight $a$, marking a
$3$ with $j$ letters to its right replaces its block weight by
$z^3a^{k-1}t^j$, for $0\le j<k$. Multiply by $t$ to the total length of
subsequent blocks. In the safe forms above, marking the $V^k$ blocks gives
the second term of~\eqref{cmb:marked-safe}; marking the first $U$ blocks
gives the third. In the early form, mark in order the $U^k$ block, the
second $U^h$ block, and the first $U^h$ block. These give the three
nonterminal terms of~\eqref{cmb:marked-early}, respectively. Summing each
block by the finite geometric identity
$\sum_{j=0}^{k-1}t^j=(1-t^k)/(1-t)$ yields the displayed expressions.
Thus these nonnegative series mark actual occurrences, with multiplicity
when a prefix has several $3$'s.

For a mark at distance $d$, let $Q_d(z)$ count binary tails that fail to
witness it. Then
\begin{equation}\label{cmb:failure-tail}
 Q_d(z)=\sum_{0\le\ell<d}c^\ell+c^dA_w(z),\qquad
 A_w(z)=\frac{1+z^2}{1-2z^3-z^4}.
\end{equation}
Indeed no witness is possible for $\ell<d$. For $\ell\ge d$, the first
$\ell-d$ positions are paired by reflection. A pair cannot be $11$, so its
possibilities are $12,21,22$, of weight $2z^3+z^4$; a central singleton
must be $2$. The last $d$ positions are arbitrary. The empty constrained
block when $\ell=d$ is included by the numerator of $A_w$.

For a marked series $F(z,t)=\sum_{d\ge0}F_d(z)t^d$, define
\begin{equation}\label{cmb:failure-operator}
 \mathcal U(F)=z\left\{\frac{F(z,1)-F(z,c)}{1-c}+A_wF(z,c)\right\}
             =z\sum_{d\ge0}F_d(z)Q_d(z).
\end{equation}
This is a nonnegative formal series counting failed marks. It upper bounds
bad bases, since each bad base has at least one failed mark. The extra
factor $z$ converts base weight into the source genus in
Lemma~\ref{cmb:witness-unused}. Set
\begin{align}
 T_s&=\frac{zC_s}{1-c},&T_e&=\frac{zC_e}{1-c},&
 U_s&=\mathcal U(F_s),&U_e&=\mathcal U(F_e).\nonumber
\end{align}
\begin{equation}
 L_{\rm nonbin}(g)=\max\{0,[z^g](T_s-U_s)\}
                  +\max\{0,[z^g](T_e-U_e)\}.\label{cmb:nonbinary-capacity}
\end{equation}
Separate clipping is legitimate because the two base families are disjoint.
Lemma~\ref{cmb:witness-unused} proves that this is a lower bound on the
number of unused nonbinary targets.

\begin{lemma}\label{cmb:nonbinary-numerical}
For $g\ge305$, $L_{\rm nonbin}(g)\ge(927/500)\Fib_g$.
\end{lemma}
\begin{proof}
Define nonnegative error series
\[
 E_s=\frac{(1+z)(1+z^2)^2}{1-v},\qquad
 E_e=\frac{z^4(1+z+z^2)}{(1-u)(1-u^2)},\qquad S=E_s+E_e+U_s+U_e.
\]
Direct multiplication of denominators gives
\begin{align*}
 T_s&=\frac1{1-c}-E_s,\\
 T_e&=\frac{z^4}{1-c}-E_e+
       \frac{z^9(2+2z+z^2)}{(1-u)(1-u^2)}.
\end{align*}
As $[z^j](1-z-z^2)^{-1}=\Fib_{j+1}$, the sum before clipping is at least
$\Fib_{g+1}+\Fib_{g-3}-[z^g]S$ for $g\ge4$.
All geometric denominators in $E_s,E_e,A_w,F_s(z,c),F_e(z,c)$ are
positive at $z=5/8$. Although $c(5/8)>1$, the expression for $\mathcal U$ is convergent
there: for this positive value of $c$,
$\sum_{\ell<d}c^\ell\le c^d/(c-1)$, so~\eqref{cmb:failure-operator}
is dominated by a constant times $F(z,c)$. We are not expanding
$(1-c)^{-1}$ as a convergent geometric series at this point.

The exact rational inputs, computed by~\eqref{cmb:coefficient-recurrence}, are
\begin{align*}
 S(5/8)&<78,\\
 78(809/500)^2(800/809)^{1500}&<1/10000,\\
 10000(\Fib_{g+1}+\Fib_{g-3})&\ge18541\Fib_g &&(g=16,17),\\
 500L_{\rm nonbin}(g)&\ge927\Fib_g &&(305\le g\le1499).
\end{align*}
The difference in the third line satisfies the Fibonacci recurrence, so
the inequality holds for every $g\ge16$. Nonnegative coefficients and
$\Fib_g\ge(809/500)^{g-2}$ give
\[
 \frac{[z^g]S}{\Fib_g}<78(809/500)^2(800/809)^g<\frac1{10000}
 \qquad(g\ge1500).
\]
Subtracting this error gives $18540\Fib_g/10000=927\Fib_g/500$ for
$g\ge1500$. The finite interval in the last line completes the proof.
\end{proof}

\begin{proposition}\label{cmb:capacity}
Define the exact integer lower bound
\begin{equation}\label{cmb:L}
 L(g)=L_{\rm bin}(g)+L_{\rm nonbin}(g).
\end{equation}
At least $L(g)$ low-depth targets of genus $g+1$ are outside the image of
$\Phi_g$. For $g\ge305$,
\begin{align}
 L(g)&\ge\frac{2841}{1000}\Fib_g,\label{cmb:capacity-constant}\\
 \frac{L(g)}{\Fib_g}&\ge1-
 \frac{50(809/500)^2}{((63/100)(809/500))^g}+\frac{927}{500}.
 \label{cmb:capacity-improved}
\end{align}
\end{proposition}
\begin{proof}
The binary and nonbinary unused sets are disjoint. Thus their bounds add.
The first estimate combines $987/1000$ and $927/500$; the second combines
$L_{\rm bin}(g)\ge\Fib_g-W_g$ with~\eqref{cmb:binary-decay} and
Lemma~\ref{cmb:nonbinary-numerical}.
\end{proof}

\section{The remaining counting problem}
\label{cmb:envelope}
Let $t_g$ count low-depth genus-$g$ words, and write
$c_g=|\mathcal C_g|$, $f_g=|\mathcal T_g|$,
$d_g=|\mathcal D_g|$, $i_g=|\mathcal I_g|$, $b_g=|\mathcal B_g|$.
Only the ordinary-child map uses high-depth targets. The exact number of
unused low-depth targets is therefore
\begin{equation}\label{cmb:exact-capacity}
 R_g=t_{g+1}-t_g-c_g-f_g-d_g-i_g-b_g\ge L(g).
\end{equation}
Counting the $n_g-|\mathcal E_g|$ matched targets and these unused targets
proves the fundamental inequality
\begin{equation}\label{cmb:fundamental}
 n_{g+1}-n_g\ge R_g-|\mathcal E_g|\ge L(g)-|\mathcal E_g|.
\end{equation}
There may also be unused high-depth targets, which only strengthen it.

Let $A_g$ count all genuine depth-four words with at least two $4$'s, and
let $B_g$ count all genuine depth-five words with at least two $5$'s.
No leaf or $1$ condition is imposed in these two definitions.

\begin{lemma}\label{cmb:parent-reduction}
The unique-$5$ sector of $\mathcal E_g$ injects into nonleaf depth-four
words of genus $g-1$ with at least two $4$'s and at least one $1$.
Consequently $|\mathcal E_g|\le A_g+A_{g-1}+B_g$.
\end{lemma}
\begin{proof}
Restore the Frobenius, changing the unique $5$ to $4$. The result is its
ordinary parent, has one new and at least one preexisting $4$, and retains
its $1$'s. It is a nonleaf because the source is its child. Two such sources
cannot have the same parent by Lemma~\ref{cmb:parent}. The other two
sectors of Lemma~\ref{cmb:residual} are directly counted by $A_g$ and $B_g$.
\end{proof}

For $q\in\{4,5\}$ define the \emph{raw cap-$q$ population} as follows.
Choose a word with at least two $q$'s and denote their two rightmost
positions by $r<p$. Coordinates before $r$ belong to $\{1,\ldots,q\}$;
all nonpivot coordinates after $r$ belong to $\{1,\ldots,q-1\}$.
Retain only the Kunz inequalities targeting $r$ or $p$, including all
self-pairs. Every genuine word counted by $A_g$ or $B_g$ belongs to the
corresponding raw population. The two rightmost maxima are unique, so
summing over their positions introduces no multiple counting.

\begin{proposition}\label{cmb:raw-envelope}
Let $U_g$ be any integer upper bound on the complete raw cap-four
coefficient of weight $g$, and let $V_g$ be any integer upper bound on
$B_g$. With $N_g=[z^g]N(z)$ from~\eqref{cmb:N}, the quantity
\begin{equation}\label{cmb:Eupper}
 E_{\rm upper}(g)=U_g-N_g-N_{g-3}-N_{g-4}
                         +U_{g-1}-N_{g-1}+V_g
\end{equation}
is an upper bound on $|\mathcal E_g|$.
\end{proposition}
\begin{proof}
Raw cap-four words without $1$ are exactly the $\{2,3,4\}$ words with
at least two $4$'s, and every such word is genuine: noncarry sums are at
least four and carry sums at least five. Their exact series is $N(z)$.
For a word whose sole $1$ is final, all inequalities are automatic except
its self-pair, which is a carry pair targeting the penultimate coordinate
and forces that coordinate to be at most three. If the penultimate
coordinate were $4$, it would be the rightmost maximum, so the raw
population retains precisely this obstruction. Thus raw and genuine
populations of this description coincide and consist of
$(v,2,1)$ and $(v,3,1)$ with $v$ a $\{2,3,4\}$ word containing at least
two $4$'s. Their series is $(z^3+z^4)N(z)$, disjoint from the no-$1$
population. Subtracting these exact counts from $U_g$ bounds the
first residual sector.

Lemma~\ref{cmb:parent-reduction} bounds the second sector by the
raw cap-four count of weight $g-1$ containing $1$, hence by
$U_{g-1}-N_{g-1}$. Its image need not be a residual leaf and may have its
sole $1$ final, so there is no further subtraction from this shifted term.
Finally $V_g$ bounds the third sector. Adding these bounds proves the
claim. The exact exclusions are from the whole raw population; the way an
upper bound $U_g$ is assembled does not affect this subtraction.
\end{proof}

For example, $(5,4,2,1)$ is a leaf with parent $(4,4,2,1)$, whose sole $1$
is final. The latter has effective generator $21$, and its deletion gives
the former. Thus the absence of a final-$1$ exclusion in the shifted term
of~\eqref{cmb:Eupper} is essential.

We will also use the following real-variable improvement. Let $D(z)$ count
the first two residual sectors together, with their actual genus weights,
and put $A(z)=\sum_gA_gz^g$, $B(z)=\sum_gB_gz^g$.
\begin{lemma}\label{cmb:compensation}
At every $0<z\le1$ for which $A(z)$ converges, $D(z)\le A(z)$.
In particular, whenever $0<z_4,z_5\le1$ and both masses are finite,
\begin{equation}\label{cmb:mass-extraction}
 |\mathcal E_g|\le A(z_4)z_4^{-g}+B(z_5)z_5^{-g}.
\end{equation}
\end{lemma}
\begin{proof}
Map each depth-four residual leaf to itself, and each unique-$5$ residual
leaf to its parent as in Lemma~\ref{cmb:parent-reduction}. The former images
are leaves and the latter are nonleaves, so they are disjoint. This gives
an injection into the population counted by $A$. The first map preserves
genus, while the second lowers it by one. For $0<z\le1$ each source weight
is therefore at most its image weight. Nonnegative summation proves
$D(z)\le A(z)$. Extract a single nonnegative coefficient at $z_4$, and bound
the remaining multiple-$5$ sector using $B$ at $z_5$, to obtain the result.
This is a positive-real comparison, not a coefficientwise assertion
$[z^g]D\le A_g$.
\end{proof}

\section{A coefficient bound for the raw two-pivot tail}\label{raw:section}

This section concerns a relaxation of the Kunz inequalities.  Its objects need
not represent numerical semigroups.  Transfer matrices and rational generating
series are established tools in the enumeration of numerical semigroups; see
Bacher~\cite{Bacher}.  The point here is a uniform coefficient estimate after
summing over both the multiplicity and an unbounded set of pivot positions.

\begin{definition}[Raw two-pivot polynomial]\label{raw:definition}
For integers $1\le r<p<m$ and $a\in\{4,5\}$, let
$Z^{(a)}_{m,r,p}(z)$ count words $(x_1,\ldots,x_{m-1})$ with
\[
 x_r=x_p=a,\qquad
 1\le x_i\le a\ (i<r),\qquad
 1\le x_i\le a-1\ (i>r,\ i\ne p),
\]
weighted by $z^{\sum_i x_i}$.  Retain exactly the inequalities
\[
 x_i+x_j\ge a\quad(i+j\in\{r,p\}),\qquad
 x_i+x_j\ge a-1\quad(i+j\in\{m+r,m+p\}),
\]
including the cases $i=j$.  No inequality with another target is imposed.
Write $Z_{m,r,p}=Z^{(4)}_{m,r,p}$ and
\[
 F_{r,p}(z)=\sum_{m>p}Z_{m,r,p}(z).
\]
\end{definition}

The pivots are uniquely the two rightmost occurrences of the maximum.
Consequently summing over $m,r,p$ counts each raw word once.  A genuine
depth-$a$ Kunz word with at least two maxima belongs to this relaxation.

\begin{theorem}[Raw omitted-pivot estimate]\label{raw:tail}
For every integer $g\ge0$,
\begin{equation}\label{raw:main}
 \sum_{p>300}\sum_{r=1}^{p-1}[z^g]F_{r,p}(z)
 \le \frac{221715}{1000000}\,\Fib_g.
\end{equation}
The left side vanishes for $g<604$.  No leaf condition or occurrence of a
coordinate equal to one is required.
\end{theorem}

\subsection{The reflection graph and its support}

Remove the two fixed pivots.  A noncarry pair gives a \emph{strong} edge,
of threshold $a$, and a carry pair gives a \emph{carry} edge, of threshold
$a-1$.  A self-pair is a unary restriction, represented as a loop rather
than an ordinary edge.  Each vertex has at most one reflection incidence
for each target.  Thus the ordinary graph has maximum degree two and its
components are paths, isolated vertices, and cycles.  A cycle alternates
the two targets modulo $m$, has even length, and has length at least four:
distinct targets cannot give parallel edges.  Constraints incident to a
fixed pivot are automatic.

This description also specifies the exact polynomial: sum the products
of vertex weights over allowed assignments along each path; for a cycle,
fix one starting value and close the last edge back to that value.
Multiply the component polynomials and the pivot factor $z^{2a}$.
Loops impose their unary restrictions at the corresponding endpoints.
An isolated vertex is a path of zero ordinary-edge length.

\begin{lemma}[Finite support]\label{raw:support}
Every word counted by $Z^{(a)}_{m,r,p}$ has genus satisfying
\begin{equation}\label{raw:genus-support}
 2g\ge(a-1)m+p+1.
\end{equation}
In particular, for cap four, $g\ge2p+2$ and $g\ge m+5$.
\end{lemma}
\begin{proof}
The reflection targeting $p$ pairs the $p-1$ coordinates to its left
with sum at least $a$ and the $m-p-1$ coordinates to its right with sum
at least $a-1$.  A central fixed point satisfies twice its coordinate
at least the same threshold.  Adding these inequalities, with twice the
pivot weight, gives
\[
 2g\ge a(p-1)+(a-1)(m-p-1)+2a=(a-1)m+p+1.
\]
For $a=4$ use $m\ge p+1$.  Independently, the two pivot entries contribute
eight and the other $m-3$ coordinates contribute at least one each,
giving $g\ge m+5$.
\end{proof}

All series and coefficient sums below are therefore well defined formally.
The support assertion in Theorem~\ref{raw:tail} follows already from
$p\ge301$.

\subsection{A structural recurrence and cancellation of pure modes}\label{raw:topology}

Fix $r<p$, and put $d=p-r$.  For $m>2p$, temporarily include the pivot
$p$ as an ordinary vertex of the carry graph.  Set $x=i-r$ and $N=m-r$.
The carry vertices are $1,\ldots,N-1$, and their two reflections are
\begin{equation}\label{raw:reflections}
 \mathcal F(x)=N-x,\qquad \mathcal G(x)=N+d-x.
\end{equation}
The first is defined throughout this interval.  The second stays in the
interval exactly when $x>d$.  Thus the vertices $1,\ldots,d$ have one
carry incidence each and form the \emph{ports}.  The pivot $p$ is port
$d$; the surviving ports correspond to coordinates $r+1,\ldots,p-1$.
Their alphabets are all $\{1,2,3\}$.

Here is the complete path decomposition, including the loops.  For a
representative $a\in\{1,\ldots,d\}$, list its residue class in the interval as
$a_k=a+kd$, $0\le k<n_a$.  Let $b\in\{1,\ldots,d\}$ be determined by
$a+b\equiv N\pmod d$.  The map $\mathcal F$ reverses the two lists and
bijects them, so $n_a=n_b=n$.  Since $a+(n-1)d+b=N$, their index formulas are
\[
 \mathcal F(a_k)=b_{n-1-k},\qquad
 \mathcal G(a_k)=b_{n-k}\quad(1\le k<n).
\]
For $a\ne b$, alternating these formulas gives the path
\[
 a_0,b_{n-1},a_1,b_{n-2},\ldots,a_{n-1},b_0.
\]
It contains all $2n$ vertices in the two classes, has ports $a,b$, and
has $2n-1$ ordinary edges.  For $a=b$, the same alternation visits the
indices $0,n-1,1,n-2,\ldots$ until all $n$ vertices have appeared.  It
has one port, $n-1$ ordinary edges, and a loop at its other endpoint.
That loop belongs to $\mathcal F$ when $n$ is odd and to $\mathcal G$
when $n$ is even.  Both loops require the coordinate value to be at
least two.  The orbits of $a\mapsto N-a\pmod d$ exhaust the residue
classes, so this list omits no carry component.

Replacing $m$ by $m+d$ leaves the port involution unchanged and adds one
vertex to every residue class.  Each paired-class path gains two
ordinary edges, and each fixed-class path gains one.  If $h_c$ is the
number of paired-class paths and $f$ the number of fixed classes, then
\[
 2h_c+f=d,
\]
so the total carry-edge increment is exactly $d$.  A change of loop
target does not change its unary restriction.  Deleting port $d$ removes
an endpoint and its incident ordinary edge, never an internal vertex.
It creates one free endpoint.  The strong graph, its loops and the
alphabets on all surviving ports are independent of $m$.

The symmetric cap-three carry matrix is
\[
 C(z)_{ij}=z^{(i+j)/2}{\bf1}_{i+j\ge3}\qquad(1\le i,j\le3).
\]
It has rank two.  Its nonzero eigenvalues $\lambda,\mu$ satisfy
\begin{equation}\label{raw:eigenvalues}
 \lambda+\mu=b=z^2+z^3,\qquad
 \lambda\mu=c=-z^3-z^4.
\end{equation}
For a positive ordinary-edge length $\ell$, its path kernel is
$C^\ell=\lambda^\ell P_\lambda+\mu^\ell P_\mu$.
There is one possible zero-edge case to check.  Since $m>2p$, every
residue class has at least two vertices.  Consequently a paired-class
path still has positive length after deleting $p$, whereas a fixed-class
path with two vertices becomes one free, loop-constrained vertex.
The free and loop boundary vectors for this component are
\[
 h=(z^{1/2},z,z^{3/2})^T,\qquad h_+=(0,z,z^{3/2})^T.
\]
They both belong to the image of $C$, which is the span of
$(1,0,0)^T$ and $(0,z,z^{3/2})^T$.  Therefore
\[
 h^T(P_\lambda+P_\mu)h_+=h^Th_+.
\]
The zero-edge component is thus also exactly the sum of its two
nonzero modes.  It introduces no extra transient term into the
recurrence.

\begin{lemma}[Fixed-pivot rationality]\label{raw:rationality}
On each progression $m=m_0+kd$ with $m_0>2p$, the polynomials
$Z_{m,r,p}$ satisfy the scalar recurrence with characteristic polynomial
\[
 P_d(T,z)=\prod_{j=0}^d(T-\lambda^{d-j}\mu^j)\in\mathbb Z[z,T].
\]
Consequently $F_{r,p}$ is a rational function.
\end{lemma}
\begin{proof}
On this progression the port pairing, endpoint restrictions and strong
core are fixed.  Expand each carry kernel into its two spectral
summands and contract against that core.  A paired-class component
assigned the $\lambda$ mode contributes a factor $\lambda^{2k}$ under
the shift $m_0\mapsto m_0+kd$; a fixed-class component contributes
$\lambda^k$.  The same statements hold with $\mu$ in place of $\lambda$.
If $j$ is the sum of the increments assigned to $\mu$, the remaining
increments sum to $d-j$.  Collecting the finitely many mode assignments
therefore gives
\[
 Z_{m_0+kd,r,p}=\sum_{j=0}^d A_j(z)(\lambda^{d-j}\mu^j)^k.
\]
This identity holds over the corresponding algebraic function field,
including at $k=0$ by the zero-edge calculation above; its amplitudes
need not be positive.  Applying $\prod_{j=0}^d(T-\lambda^{d-j}\mu^j)$
to this exponential sum proves the recurrence.
For integrality of the annihilator, put
$S_0=2$, $S_1=b$, and $S_n=bS_{n-1}-cS_{n-2}$.
Pair the roots indexed by $j,d-j$ to obtain
\[
 T^2-c^jS_{d-2j}T+c^d\qquad(0\le j<d/2),
\]
with the additional factor $T-c^{d/2}$ when $d$ is even.

For an explicit rational construction define
\[
 Q_d(y,z)=\prod_{j=0}^d(1-\lambda^{d-j}\mu^j y)
          =\sum_{j=0}^{d+1}q_j(z)y^j.
\]
For $1\le s\le d$ set $m_s=2p+s$ and
\[
 N_s(y,z)=\sum_{k=0}^d\sum_{j=0}^k
              q_j(z)Z_{m_s+(k-j)d,r,p}(z)y^k.
\]
The recurrence yields
\begin{equation}\label{raw:initial-data}
 F_{r,p}(z)=\sum_{m=p+1}^{2p}Z_{m,r,p}(z)
              +\frac{\sum_{s=1}^dN_s(1,z)}{Q_d(1,z)}.
\end{equation}
The substitution $y=1$ is justified coefficientwise by $g\ge m+5$,
and $Q_d(1,0)=1$.  In particular this is a structural recurrence,
not a recurrence fitted to finitely many values.
\end{proof}

\begin{proposition}[Cancellation and numerator degree]\label{raw:cancellation}
Let $\tau=1$ when $d$ is odd and $\tau=2$ when $d$ is even.  Then
\begin{equation}\label{raw:correction}
 H_{r,p}:=(1-z-z^2)F_{r,p}=\frac{A_{r,p}}{B_d},\qquad
 B_d=K_\tau R_d,
\end{equation}
where $A_{r,p}\in\mathbb Q[z]$, $\deg A_{r,p}\le3d^2+7p+6$, and
\begin{align*}
 K_1&=1+z+z^2,\\
 K_2&=1+z+2z^2+z^3-z^5-z^6,\\
 R_d&=\prod_{j=1}^{d-1}(1-\lambda^{d-j}\mu^j)\in\mathbb Z[z].
\end{align*}
\end{proposition}
\begin{proof}
We first separate the multiplicity dependence of a pure mode from the
way in which its ports are paired.  Fix $\theta\in\{\lambda,\mu\}$,
choose a nonzero eigenvector $v_\theta$ of the symmetric carry matrix,
and put $t_\theta=v_\theta^Tv_\theta$.  Over an algebraic function field
these choices can be made with $t_\theta\ne0$, and
\[
 P_\theta=\frac{v_\theta v_\theta^T}{t_\theta}.
\]
Let $\mathcal C_{r,p}(v_\theta)$ be the contraction of the fixed strong
core, including its strong loops and the factor $z^8$ for the pivots.
At each surviving port with coordinate value $i$, use the weight
$z^{i/2}(v_\theta)_i$; use the full weight $z^i$ at each other core
vertex.  The other half of each port weight is already in the adjacent
symmetric carry kernel.  This definition specifies one fixed scalar,
independent of $m$ and of the port pairing.

There are $d-1$ surviving ports, one free carry endpoint created by
deleting $p$, and $f$ carry-loop endpoints.  There are
$(d+f)/2$ carry components, each contributing one rank-one
projector.  Factoring its two endpoint entries shows that the entire
pure-$\theta$ contribution is exactly
\begin{equation}\label{raw:pure-contraction}
 \theta^{e_C(m)}\mathcal C_{r,p}(v_\theta)
 \frac{(h^Tv_\theta)(h_+^Tv_\theta)^f}
      {t_\theta^{(d+f)/2}}.
\end{equation}
For example, a paired-port kernel contributes
$\theta^\ell(v_\theta)_i(v_\theta)_j/t_\theta$; a port-to-loop
kernel replaces one endpoint entry by $h_+^Tv_\theta$.  The component
containing the deleted port instead has the factor $h^Tv_\theta$ at
its free endpoint.  The zero-edge free-to-loop component has these
same factors, by its full two-projector decomposition in
Section~\ref{raw:topology}.  Thus every surviving port contributes
exactly one entry of $v_\theta$, regardless of its partner.
Formula~\eqref{raw:pure-contraction} is unchanged if $v_\theta$ is
rescaled, and proves the asserted independence from pairing.

For $m>2p$, counting the two carry reflections and deleting the
ordinary edge at $p$ gives
\[
 e_C(m)=\left\lfloor\frac{m-r-1}{2}\right\rfloor+
         \left\lfloor\frac{m-p-1}{2}\right\rfloor-1.
\]
The fixed classes solve $2a\equiv m-r\pmod d$.  When $d$ is odd,
there is one such class, so $f=1$ and
$e_C=m-(r+p+5)/2$.  All factors in
\eqref{raw:pure-contraction} other than $\theta^{e_C(m)}$ are
therefore constant as $m$ increases by one.  Summing the pure terms
has denominator $1-\theta$.  When $d$ is even, the values are
\[
 \begin{array}{c|cc}
 &f&e_C(m)\\\hline
 m-r\text{ odd}&0&m-(r+p)/2-2\\
 m-r\text{ even}&2&m-(r+p)/2-3
 \end{array}
\]
and splitting the sum by parity gives denominator $1-\theta^2$.
In either case the two pure contributions together have a denominator
dividing $E_\tau$, where
\[
 E_j=(1-\lambda^j)(1-\mu^j)=1-S_j+c^j,\qquad
 L_d=E_d/E_\tau,\qquad \Delta=b^2-4c.
\]
Since $\tau\mid d$, the quotient $L_d$ is a polynomial: it is the
product of the two finite geometric sums for $\lambda^\tau$ and
$\mu^\tau$, which is symmetric in $\lambda,\mu$.

We next descend this mode calculation to rational functions with
controlled denominators.  This avoids assuming that the roots being
cancelled are simple.  The rational-weight carry matrix is
\[
 M=\begin{pmatrix}0&z^2&z^3\\z&z^2&z^3\\z&z^2&z^3\end{pmatrix}.
\]
It is similar to $C$ and assigns the full vertex weight to the arriving
vertex of an oriented edge.  Its nonzero spectral projector is
\[
 \Pi_\lambda=\frac{M(M-\mu I)}{\lambda(\lambda-\mu)},
\]
with the analogous formula for $\Pi_\mu$.  Set
\[
 \mathcal A=\mathbb Q[z,(c\Delta)^{-1}],\qquad
 \mathcal S=\mathcal A[\lambda].
\]
The identities $\lambda^{-1}=\mu/c$ and
$(\lambda-\mu)^{-1}=(\lambda-\mu)/\Delta$ put both projectors in
$\operatorname{Mat}_3(\mathcal S)$.  All strong-core and boundary
weights in this rational convention also belong to $\mathcal S$;
any redistribution of a port weight uses only powers of $z$, which
are invertible in $\mathcal A$.  Thus finite mode amplitudes have no
possible denominators outside this localization.

The polynomial
$\Delta=z^3(z+1)(z^2+z+4)$ is not a square in $\mathbb Q(z)$, so
\[
 \mathcal S=\mathcal A\oplus\mathcal A\lambda.
\]
The conjugation $\sigma(\lambda)=\mu=b-\lambda$ has fixed ring
$\mathcal S^\sigma=\mathcal A$: if $a_0+a_1\lambda$ is fixed, then
$a_1(\lambda-\mu)=0$, hence $a_1=0$.
For the pure terms, the preceding summation has denominator dividing
$E_\tau$.  For a mixed assignment on any multiplicity progression,
Lemma~\ref{raw:rationality} gives a geometric denominator
$1-\lambda^{d-j}\mu^j$ with $1\le j\le d-1$, which is a factor of
$R_d$.  It follows that $E_\tau R_dF_{r,p}$ belongs to $\mathcal S$.
The automorphism $\sigma$ of $\mathbb Q(z)(\lambda)$ fixes
$\mathbb Q(z)$ pointwise. By Lemma~\ref{raw:rationality},
$F_{r,p}\in\mathbb Q(z)$, while $E_\tau,R_d\in\mathbb Q[z]$;
hence their product is fixed by $\sigma$. Therefore
\begin{equation}\label{raw:localization}
 E_\tau R_dF_{r,p}\in\mathcal A.
\end{equation}

On the other hand, \eqref{raw:initial-data} writes
$F_{r,p}=N/(E_dR_d)$ with $N\in\mathbb Q[z]$.
Equation~\eqref{raw:localization} says precisely that $N/L_d$ belongs
to $\mathcal A$.  To conclude polynomial divisibility, we check
$\gcd(L_d,c\Delta)=1$.  At $z=0,-1$ both eigenvalues vanish and
$E_d=1$.  At the other discriminant zeros, $z^2+z+4=0$, the eigenvalues
coincide at $b/2=-2z$ and have modulus four, so again $E_d\ne0$.
Since $L_d$ divides $E_d$, it has no zero at any zero of $c\Delta$.
Write $N/L_d=G/(c\Delta)^k$ in \eqref{raw:localization}.  Then
$N(c\Delta)^k=GL_d$, and coprimality implies $L_d\mid N$ in
$\mathbb Q[z]$, with every root multiplicity.  This argument does not
require $L_d$ and $R_d$ to be coprime.  Finally
\[
 E_\tau=(1-z-z^2)K_\tau
\]
gives \eqref{raw:correction} with $A_{r,p}=N/L_d$.

For the degree bound, the maximum allowed weight of a graph word is
\[
 8+4(r-1)+3(m-r-2)=3m+r-2.
\]
In the paired-factor formula for $Q_d(y,z)$, the coefficients of
$y$ and $y^2$ have degrees at most $3d$ and $4d\le6d$, respectively;
the optional middle factor has $y$-coefficient of degree $2d\le3d$.
Consequently $\deg q_j\le3dj$.  A summand in $N_s$ has degree at most
\[
 3dj+3\bigl(m_s+(k-j)d\bigr)+r-2
 \le D_0:=3d^2+3d+6p+r-2,
\]
since $m_s\le2p+d$ and $k\le d$.
Each paired factor of $Q_d(1,z)$ has degree $4d$, and the optional
middle factor has degree $2d$, giving
$\deg Q_d(1,z)=2d(d+1)$.  Thus the finite prefix in
\eqref{raw:initial-data}, after multiplication by this denominator,
contributes degree at most
\[
 2d(d+1)+6p+r-2\le D_0.
\]
Finally $\deg E_j=4j$, so cancelling $L_d$ removes degree
$4(d-\tau)$.  Substituting $r=p-d$ yields
\[
 \deg A_{r,p}\le3d^2-2d+7p+4\tau-2\le3d^2+7p+6.
\]
\end{proof}

\subsection{Uniform absolute coefficient norms}

Set $q=(\sqrt5-1)/2$, and for a formal series $f$ write
$\|f\|_q=\sum_{j\ge0}|[z^j]f|q^j$, allowing the value infinity.
This norm is submultiplicative whenever the right side is finite.

\begin{lemma}[Uniform denominator bounds]\label{raw:wiener}
For every $d\ge1$,
\begin{equation}\label{raw:norms}
 \|B_d\|_q<10,\qquad \|B_d^{-1}\|_q<21.
\end{equation}
\end{lemma}
\begin{proof}
At $q$ the nonzero carry eigenvalues are $1,-q^2$.  Put $u=q^2$.
The polynomials $S_n$ have nonnegative coefficients, since $b,-c$ do.
For a paired mixed factor, $j\ge1$ and $n=d-2j\ge1$, its nonconstant
coefficient norm is at most
\[
 u^jS_n(q)+u^d\le u^j+u^{d-j}+u^d
                  \le a_j:=u^j+u^{j+1}+u^{2j+1}.
\]
Let
\[
 a_1=u+u^2+u^3,\qquad
 T=\sum_{j\ge2}a_j=\frac{(1+u)u^2}{1-u}+\frac{u^5}{1-u^2}.
\]
Both $a_1,T$ are less than one.  An optional central factor
$1-c^{d/2}$ costs at most $1+u$ in norm and $(1-u)^{-1}$ in
inverse norm.  Geometric-series bounds and the finite-product inequality
$\prod(1-a_j)\ge1-\sum a_j$ give
\begin{equation}\label{raw:mixed-norms}
 \|R_d\|_q\le\frac{(1+u)(1+a_1)}{1-T}<\frac{10}{3},\qquad
 \|R_d^{-1}\|_q\le\frac1{(1-u)(1-a_1)(1-T)}<6.
\end{equation}
These comparisons follow on replacing $q$ by the larger rational number
$q_+=309017/500000$, since the displayed upper bounds increase with $u$.

For completeness the fixed-factor reciprocal estimates can be certified
by a short, precisely specified finite calculation.  For either coefficient
list
\[
 (k_0,\ldots,k_e)=(1,1,1)\quad\hbox{or}\quad
 (1,1,2,1,0,-1,-1),
\]
put $v_0=1$, $v_n=-\sum_{j=1}^{\min(e,n)}k_jv_{n-j}$ for
$1\le n\le80$, $V(z)=\sum_{n=0}^{80}v_nz^n$, and $E=1-KV$.
The deliberately rounded rational inequalities in Table~\ref{tab:raw-reciprocal-residuals} are sufficient:

\begin{table}[ht]
\centering
\caption{Rational bounds for the degree-$80$ reciprocal polynomials
and their residuals at $q_+=309017/500000$.}
\label{tab:raw-reciprocal-residuals}
\begin{tabular}{c|cc}
 & $\|V\|_{q_+}$ & $\|E\|_{q_+}$\\\hline
 $K_1$ & $<212/100$ & $<1/10^{16}$\\
 $K_2$ & $<343/100$ & $<1/10^{8}$
\end{tabular}
\end{table}

They are finite rational inequalities for the recurrence just displayed:
all coefficients of $E$ through degree $80$ vanish and only the next
$e$ coefficients enter its norm.  In particular $\|E\|_{q_+}<1$,
so the identity $K^{-1}=V\sum_{j\ge0}E^j$ proves, including its
infinite tail,
\[
 \|K_1^{-1}\|_q<7/2,\qquad \|K_2^{-1}\|_q<7/2.
\]
Also $\|K_1\|_q=2$ and $\|K_2\|_q=4-2q<14/5$.
Combining these with \eqref{raw:mixed-norms} proves \eqref{raw:norms}.
\end{proof}

\begin{lemma}[A finite numerator bound and signed conversion]\label{raw:signed}
Suppose $Z_{m,r,p}(q)\le B(r,p)$ for every $m>p$.  Then, with
$N=3d^2+7p+6$,
\begin{equation}\label{raw:numerator-norm}
 \|H_{r,p}\|_q\le420N B(r,p),
\end{equation}
and, for every $g\ge0$,
\begin{equation}\label{raw:fib-conversion}
 [z^g]F_{r,p}\le\frac94\|H_{r,p}\|_q\Fib_g.
\end{equation}
\end{lemma}
\begin{proof}
Only coefficients through degree $N$ of $F_{r,p}$ can enter the polynomial
$A_{r,p}=B_d(1-z-z^2)F_{r,p}$.  Positivity of each graph polynomial
and $g\ge m+5$ imply
\[
 \sum_{j=0}^N[z^j]F_{r,p}\,q^j\le N B(r,p),
\]
since fewer than $N$ multiplicities contribute.  Using
$\|1-z-z^2\|_q=2$ and \eqref{raw:norms} gives
$\|A_{r,p}\|_q\le20NB(r,p)$ and then
\eqref{raw:numerator-norm}.  This step never substitutes the potentially
divergent value $F_{r,p}(q)$ for its finite truncation.

Write $H_{r,p}=\sum_jh_jz^j$.  Its support starts at degree at least
eight.  From $(1-z-z^2)^{-1}=\sum_{n\ge0}\Fib_{n+1}z^n$,
\[
 [z^g]F_{r,p}=\sum_{j=8}^g h_j\Fib_{g-j+1}.
\]
For $8\le j\le g$, Binet's formula and
$\Fib_n\le q^{1-n}$ give
\[
 \frac{\Fib_{g-j+1}}{\Fib_g}
 \le \frac{\sqrt5}{1-q^{16}}q^j<\frac94q^j.
\]
For the last strict inequality it suffices to square and check
$5<[(9/4)(1-q_+^{16})]^2$.  Take absolute values of the $h_j$.
For $g<8$ the raw coefficient is zero, which also handles $g=0$.
\end{proof}

\begin{remark}\label{raw:not-positive}
The signed argument is necessary for this correction.  For $r=1,p=2$,
the word $(4,4)$ contributes at genus eight, whereas no raw word for
these pivots has genus nine.  At multiplicity four the remaining
coordinate has a carry self-pair and is at least two; multiplicities
at least five have genus at least ten.  Hence
$[z^9](1-z-z^2)F_{1,2}=-1$.
\end{remark}

\subsection{Strong runs at the Fibonacci point}

For the remainder of the proof all matrices are evaluated at $q$.
On four coordinates let
\[
 D_{ij}=q^{(i+j)/2}{\bf1}_{i+j\ge4},\qquad
 C_{ij}=q^{(i+j)/2}{\bf1}_{i+j\ge3}{\bf1}_{i,j\le3}.
\]
Let $E$ be $D$ with its fourth row and column set to zero.  Thus $E$
is the alphabet-three strong transfer.  Put
\begin{gather}\label{raw:constants}
 R=\frac{917}{1000},\quad k_0=\frac{10003}{10000},\quad
 \chi=\frac{47597}{50000},\quad t=\frac{95761}{100000},\nonumber\\
 u_*=(1,13/50)^T,\quad v_*=(823/1250,97/625)^T,\quad
 \nu=v_*^Tu_*=\frac{10918}{15625}.
\end{gather}

\begin{lemma}[Uniform transfer block]\label{raw:block}
Let $Q=(v,w)$ be an orthonormal basis for the nonzero eigenspaces of
$C$, in the order $1,\mu=-q^2$.  For all $k,n\ge1$,
\begin{equation}\label{raw:block-bound}
 \left|Q^TD^kQ\begin{pmatrix}1&0\\0&\mu^n\end{pmatrix}\right|
 \le R^{\max(k-3,0)}u_*v_*^T
\end{equation}
entrywise.  Moreover $\|D\|_2<R$, $\|E\|_2<19/25$, and
$Q^TEQ=Q^TDQ$.
\end{lemma}
\begin{proof}
The nonsymmetric matrix similar to the cap-$a$ restriction of $D$ has
entries $q^j{\bf1}_{i+j\ge4}$.  Its products with the positive vectors
$(10000,17164,23907,23907)^T$ for $a=4$ and $(1,2,3)^T$ for $a=3$
are componentwise strictly smaller than $R$ and $19/25$ times those
vectors, respectively.  Each comparison reduces to $q^2=1-q$.
The positive-vector bound on spectral radius, and symmetry of the
original matrices, give the asserted operator norms.

One may choose
\[
 v=\frac{(q,\sqrt q,q,0)^T}{\sqrt{1+q^2}},\qquad
 w=\frac{(1,-q^{3/2},-q^2,0)^T}{\sqrt{1+q^2}}.
\]
Their required moments, writing
$\alpha_k=v^TD^kv$, $o_k=v^TD^kw$, and $\beta_k=w^TD^kw$, are
\begin{align*}
 \alpha_1=\alpha_2&=(7-6q)/5,&
 o_1&=(1-3q)/5,& o_2&=(11-18q)/5,\\
 \beta_1&=(-7+11q)/5,&
 \beta_2&=(3-4q)/5,& \beta_4&=(-181+293q)/5.
\end{align*}
These exact identities give the rational bounds
\[
 \begin{gathered}
 \alpha_1,\alpha_2<A=\frac{823}{1250},\quad
 |o_1|<\frac{171}{1000},\quad |o_2|<\frac1{40},\\
 |\beta_1|<\frac{41}{1000},\quad
 \beta_2<B=\frac{66}{625},\quad \beta_4<J=\frac{17}{1000}.
\end{gathered}
\]
Also $|\mu|<\rho=191/500$.  For $k\ge3$, Cauchy--Schwarz gives
\[
 |\alpha_k|\le R^{k-2}A,\quad
 |o_k|\le R^{k-3}\sqrt{AJ},\quad
 |\beta_3|\le\sqrt{BJ},\quad
 |\beta_k|\le R^{k-4}J\quad(k\ge4).
\]
For example, the middle bound follows from
$(Dv)^TD^{k-3}(D^2w)$.
To check all powers, write $M_*=u_*v_*^T$.  For $k=1,2$ use the
listed moments, multiplying the second column by at most $\rho$.
For $k\ge3$ the sufficient rational comparisons are
\[
 RA<M_{*,11},\quad AJ<M_{*,21}^2,\quad
 \rho^2AJ<M_{*,12}^2,\quad
 \rho^2BJ<M_{*,22}^2,\quad \rho J/R<M_{*,22}.
\]
Direct substitution proves them.  These inequalities prove
\eqref{raw:block-bound} for every power, rather than only for a finite
list of lengths.  The fourth coordinates of $v,w$ vanish, so
$Q^TEQ=Q^TDQ$.
\end{proof}

\begin{proposition}[Uniform graph-value estimate]\label{raw:pointwise}
Define
\[
 e_D=\left\lfloor\frac{r-1}{2}\right\rfloor+
     \left\lfloor\frac{p-1}{2}\right\rfloor-{\bf1}_{p\ne2r},
 \qquad
 e_E=\max\left(0,\left\lfloor\frac{p-1}{2}\right\rfloor-r\right),
\]
and $s_4(z)=z+z^2+z^3+z^4$.  For every $m>p$,
\begin{equation}\label{raw:value}
 Z_{m,r,p}(q)\le k_0^3q^8s_4(q)R^{e_D+2e_E}\chi^{d-16}.
\end{equation}
\end{proposition}
\begin{proof}
Call a maximal nonempty sequence of ordinary strong edges in a component
of the reflection graph a \emph{strong run}.  We first count its edges
and boundaries, then estimate its transfer product.

There are $\lfloor(r-1)/2\rfloor$ ordinary strong pairs targeting $r$
and $\lfloor(p-1)/2\rfloor$ targeting $p$.  The only such pair incident
to a pivot is $\{r,p-r\}$, unless $p=2r$, when this is a loop at $r$.
Deleting the pivots therefore leaves exactly $e_D$ ordinary strong edges.
Among these, a pair with both endpoints greater than $r$ must target $p$
and is uniquely specified by
\[
 r<i<p-i,\qquad
 i=r+1,\ldots,\left\lfloor\frac{p-1}{2}\right\rfloor.
\]
There are exactly $e_E$ such edges.  Neither endpoint can have a strong
edge targeting $r$, so each of these is an entire one-edge strong run.
It uses the alphabet-three matrix $E$.  Every other run of $k$ edges
is bounded by $D^k$, enlarging any smaller vertex alphabet as needed.
Every carry endpoint is greater than its target and has alphabet three;
the zero fourth row and column of $C$ enforce this alphabet at each
strong/carry junction.

For the boundary count, each nonpivot coordinate has one reflection
incidence for each target.  Such an incidence is either an ordinary
edge, a loop, or an edge lost on deleting a pivot.  Pivot $r$ has only
the possible nonloop neighbor $p-r$, and pivot $p$ only the possible
nonloop neighbor $m+r-p$; the pivots are not adjacent.  Hence there
are at most two free boundary incidences.  For each target $t$, loops
can occur only at $2i=t$ or $2i=m+t$, so there are at most four loop
incidences in total.  If $P$ is the number of path components, $F$ the
number of free incidences, and $L$ the number of loop incidences, then
\[
 2P=F+L,\qquad F\le2,\qquad L\le4,\qquad P\le3.
\]
An isolated vertex counts as a zero-edge path with two boundary
incidences, so it is included in this identity.

Let $h_4=(q^{1/2},q,q^{3/2},q^2)^T$ and let $h_{4,+}$ be the same
vector with its first entry zero.  Free endpoint vectors are bounded
by $h_4$, and loop endpoint vectors by $h_{4,+}$, because either loop
threshold forces a value at least two.  Thus all boundary norms
combined cost at most
\begin{equation}\label{raw:boundary-budget}
 \|h_4\|_2^F\|h_{4,+}\|_2^L
 =s_4(q)^{F/2}\bigl(s_4(q)-q\bigr)^{L/2}\le s_4(q),
\end{equation}
since $s_4(q)>1$ and $s_4(q)-q<1$.  The same formula handles the two
restrictions on an isolated vertex through the inner product of its
two boundary vectors.

Every ordinary cycle contains both strong and carry edges.  Indeed,
target labels alternate around a cycle.  In an all-strong cycle the
two target reflections compose to a translation by $d$ or $-d$;
in an all-carry cycle they do the same.  Successive pairs of edges
use the same ordering of target labels, so iterating that nonzero
translation cannot close the cycle.

For each cycle, rotate the transfer product to alternate nonempty
strong and carry runs.  Write $T_i=D^{k_i}$ for an ordinary strong
run and $T_i=E$ for an alphabet-three strong run, and put
$\Lambda_n=\operatorname{diag}(1,\mu^n)$.  Its partition is bounded by
\[
 \operatorname{tr}\bigl(T_1C^{n_1}\cdots T_jC^{n_j}\bigr)
 =\operatorname{tr}\prod_{i=1}^j(Q^TT_iQ\Lambda_{n_i}),
 \qquad n_i\ge1.
\]
By Lemma~\ref{raw:block}, the absolute value of each block is bounded
entrywise by $a_iM_*$, where $M_*=u_*v_*^T$ and
\[
 a_i=R^{\max(k_i-3,0)}\quad(T_i=D^{k_i}),\qquad
 a_i=1\quad(T_i=E).
\]
The latter follows from $Q^TEQ=Q^TDQ$.  Since
$M_*^j=\nu^{j-1}M_*$, taking absolute values and multiplying the
nonnegative comparison matrices gives the cycle bound
\[
 \left(\prod_i a_i\right)\operatorname{tr}(M_*^j)
 =\left(\prod_i a_i\right)\nu^j.
\]
In particular, no dimension factor is charged for a cycle.

For a path, remove the strong run at either endpoint if present and
charge its operator norm.  A path consisting entirely of strong edges
has just one such run.  We call these runs \emph{unpaid} only with
respect to the additional factor $\nu/R^3$ below: their edges are
already charged.  At most two runs per path are unpaid.  Unless the
whole path has been removed, what remains begins and ends with carry
runs.  If it contains $j\ge1$ strong runs, its product after projection
into the carry image has the form
\[
 \Lambda_{n_0}(Q^TT_1Q\Lambda_{n_1})\cdots
                 (Q^TT_jQ\Lambda_{n_j}),\qquad n_i\ge1.
\]
Orthogonal projection and $\Lambda_{n_0}$ have operator norm at most
one.  Taking entrywise absolute values preserves the Euclidean norm
of each endpoint vector.  The same block comparison consequently bounds
this middle product by
\[
 \left(\prod_i a_i\right)\|M_*^j\|_2
 =\left(\prod_i a_i\right)\kappa\nu^j,
 \qquad \kappa=\frac{\|u_*\|_2\|v_*\|_2}{\nu}<k_0.
\]
If there is no remaining strong run, the carry product has norm at most
one.  This includes a zero-edge remainder; a path removed in its entirety
was already bounded by its endpoint norms and strong-edge charges.
Thus in every case one factor $k_0$ per path suffices, at most $k_0^3$
in total.  We call all strong runs in cycles and all nonendpoint strong
runs in paths \emph{paid}.

To count paid runs, consider the strong graph before deleting $r$.
On its vertices $1,\ldots,p-1$, the substitution $x=p-i$ turns the
reflections into $p-x$ and $p+d-x$, with ports $1,\ldots,d$.
The residue-class path argument in Section~\ref{raw:topology} applies
here too: it only requires each residue class to be nonempty, which
holds because $p-1\ge d$.  If $f_s$ classes are fixed by the residue
involution, then $f_s\le2$, because $2a\equiv p\pmod d$ has at most
two solutions.  The $(d-f_s)/2\ge d/2-1$ paired-class components
each contain an ordinary strong edge, even if their classes have only
one vertex each.  The deleted pivot $r$ is port $d$.  Removing that
endpoint can destroy at most one of these components and cannot split
another.  At least $d/2-2$ nonempty strong components therefore survive.
These components are precisely strong runs after carry edges are
restored.  At most $2P\le6$ runs were declared unpaid, so the total
number $j$ of paid runs satisfies
\begin{equation}\label{raw:paid-run-count}
 j\ge d/2-8.
\end{equation}
This argument involves only the strong interval and the boundary count;
in particular it applies throughout $p<m\le2p$ as well as for $m>2p$.

The constants satisfy the exact rational comparisons
\[
 \frac{\nu}{R^3}<\chi^2<1,\qquad
 \frac{\|u_*\|_2\|v_*\|_2}{\nu}<k_0,\qquad
 \frac{19}{25}<R^3,\qquad \chi<t,\qquad R<t^2.
\]
Charge $R$ per ordinary strong edge outside the $E$ runs and $R^3$
per $E$ edge.  An unpaid run is covered by its operator norm:
$\|D^k\|_2\le R^k$ or $\|E\|_2<19/25<R^3$.
For a paid run, the preceding product estimates give
\[
 R^{\max(k-3,0)}\nu\le R^k\frac{\nu}{R^3}
 \quad\hbox{for a $D$ run},\qquad
 \nu=R^3\frac{\nu}{R^3}\quad\hbox{for an $E$ run}.
\]
Every strong edge is counted exactly once, so the total strong charge is
$R^{e_D-e_E+3e_E}=R^{e_D+2e_E}$.
By \eqref{raw:paid-run-count}, the extra factor is bounded by
\[
 (\nu/R^3)^j\le\chi^{2j}\le\chi^{d-16}.
\]
When $d<16$, the last negative exponent simply weakens the bound.
Multiplying this estimate by the path factor $k_0^3$, the global
boundary bound \eqref{raw:boundary-budget}, and the pivot weight $q^8$
proves \eqref{raw:value}.
\end{proof}

\subsection{Summation over omitted pivots}

\begin{proof}[Proof of Theorem~\ref{raw:tail}]
Combining Lemma~\ref{raw:signed} with Proposition~\ref{raw:pointwise}
gives the coefficient multiplier
\[
 945(3d^2+7p+6)k_0^3q^8s_4(q)
                 R^{e_D+2e_E}\chi^{p-r-16}.
\]
Replace $d^2$ by $p^2$.  Put $J=\lfloor(p-1)/2\rfloor$ and
$a=R^2\chi/t<1$.  Since
$R^{\lfloor(r-1)/2\rfloor}\le t^{r-2}$, the region $r\le J$
is bounded, after reversing its geometric sum, by
\[
 \sum_{r\le J}R^{e_D+2e_E}\chi^{p-r}
 \le\frac{\chi^{p-J}(Rt)^J}{Rt^2(1-R^2\chi/t)}.
\]
Here replacing the pivot-deletion indicator by one only enlarges the
bound.  For $r>J$ there are no $E$ edges; summing backwards from
$r=p-1$ gives
\[
 \sum_{J<r<p}R^{e_D}\chi^{p-r}
 \le\frac{R^Jt^p}{Rt^2(t/\chi-1)}.
\]
An empty region is harmless because the displayed upper bound is positive.

For a quadratic $P(k)=ak^2+bk+c$, define
\begin{equation}\label{raw:geometric-tail}
 \mathcal T_x(K;P)=x^K\left[
 \frac{aK^2+bK+c}{1-x}
 +\frac{(2aK+b)x}{(1-x)^2}
 +\frac{ax(1+x)}{(1-x)^3}\right]
 =\sum_{k\ge K}P(k)x^k.
\end{equation}
Set
\[
 X=Rt\chi,\quad Y=Rt^2,\quad U=1-R^2\chi/t,\quad V=t/\chi-1,
\]
and
\[
 P_o(k)=12k^2+26k+16,\qquad P_e(k)=12k^2+38k+32.
\]
These are $3p^2+7p+6$ for $p=2k+1,2k+2$, respectively.
In both cases $J=k$.  Pairing these two parities, and replacing $q$
by $q_+$ in the positive prefactor, bounds the multiplier for every
pivot $p>2K$ by the following explicit rational number:
\begin{equation}\label{raw:explicit-tail}
 \mathcal B_K=
 \frac{945q_+^8s_4(q_+)k_0^3\chi^{-16}}{Rt^2}
 \left[
 \frac{\mathcal T_X(K;\chi P_o+\chi^2P_e)}{U}
 +\frac{\mathcal T_Y(K;tP_o+t^2P_e)}{V}
 \right].
\end{equation}
All entries in this expression are the specified rational constants;
$0<X<Y<1$ and $U,V>0$.
Substitution in \eqref{raw:geometric-tail} and integer cross-multiplication
give the exact inequality
\[
 \mathcal B_{150}<\frac{221714122}{10^9}
                       <\frac{221715}{10^6}.
\]
This finite arithmetic comparison bounds the whole infinite geometric
sum in \eqref{raw:explicit-tail}.  Together with the preceding structural
lemmas it proves \eqref{raw:main} for every genus.  For a fixed genus
the pivot sum itself is finite by Lemma~\ref{raw:support}; the displayed
majorant additionally converges absolutely.
\end{proof}

\begin{corollary}[Integer envelope]\label{raw:integer-envelope}
Let $C_g(300)=\sum_{2\le p\le300}\sum_{r<p}[z^g]F_{r,p}$ be the
exact coefficient of the raw cap-four prefix.  Then the full raw
cap-four coefficient is at most
\[
 C_g(300)+{\bf1}_{g\ge604}
              \left\lfloor\frac{221715\Fib_g}{1000000}\right\rfloor.
\]
\end{corollary}
\begin{proof}
The omitted coefficient is an integer, so the floor of its real upper
bound is still an upper bound.  Its support starts at genus $604$.
\end{proof}

\section{Convergent masses of the genuine residual populations}
\label{mass:section}

Let $A_g$ count genuine depth-four Kunz words of genus $g$ with at
least two coordinates equal to four, and let $B_g$ count genuine depth-five
words with at least two coordinates equal to five. Write
$A(z)=\sum_g A_gz^g$ and $B(z)=\sum_g B_gz^g$.
We prove
\begin{equation}\label{mass:values}
 A(1247/2000)<2024,\qquad B(633/1000)<113.
\end{equation}
These are inequalities at specified positive real arguments. In particular,
the first inequality does not assert convergence there of the unrestricted
cap-four relaxation. The use of weighted Kunz graphs and transfer matrices
continues methods of Zhu~\cite{Zhu} and Bacher~\cite{Bacher}; the additional
self-pair restrictions below are essential to the present explicit estimate.

\subsection{The multiple-five mass}

For $t\in\{4,5\}$ let
\[
 D_t(z)_{ij}=z^{(i+j)/2}{\bf1}_{i+j\ge t},\qquad 1\le i,j\le t,
\]
and put $s_5=\sum_{j=1}^5z^j$, $h_5=z^3+z^4+z^5$ and
$h_4=z^2+z^3+z^4$. The Schatten fourth norm of a real symmetric matrix
$M$ is $(\operatorname{tr}M^4)^{1/4}$.

\begin{lemma}\label{mass:five}
Suppose $0<a_5<b_4<1$, $\|D_5(z)\|_4\le a_5$,
$\|D_4(z)\|_4\le b_4$, $s_5\ge b_4$, $h_5\le a_5$ and
$h_4\le b_4$. Set $k=(a_5/b_4)^{1/2}$. Then
\[
 B(z)\le \frac{z^{10}(s_5/b_4)k^{-1}}
 {(1-b_4)(1-b_4k)(1-a_5)}.
\]
In particular $B(633/1000)<113$.
\end{lemma}

\begin{proof}
Fix the two rightmost fives at $r<p<m$, delete these fixed vertices,
and retain the two reflection constraints. Carry endpoints have alphabet
$\{1,2,3,4\}$; enlarging the strong-endpoint alphabet to
$\{1,2,3,4,5\}$ can only increase the partition. The ordinary graph is a
union of paths and even cycles, the latter having length at least four.
Regard an isolated vertex as a path with two endpoint incidences. Let $P$
be the number of paths, $u$ the number of free incidences, $L_5,L_4$ the
numbers of strong and carry loop incidences, and $S,C$ the numbers of
ordinary strong and carry edges. Reflection counting gives
\[
 2P=u+L_5+L_4,\quad S+C=m-3-P,\quad u\le2,
 \quad 2S+L_5\ge r+p-4.
\]
The final inequality includes $p=2r$: deleting the second pivot then
removes a strong loop rather than an ordinary strong edge.
Cauchy--Schwarz on each path bounds its endpoint factors by
$s_5^{u/2}a_5^{L_5/2}b_4^{L_4/2}$ and its edge factors by
$a_5^S b_4^C$. For a cycle of length $\ell\ge4$, Schatten H\"older and
$\|M\|_\ell\le\|M\|_4$ give the same edge bound with no dimension
factor. Alphabet projections are contractions and preserve these bounds.
Consequently the nonpivot partition is at most
\[
 b_4^{m-3}(s_5/b_4)^{u/2}k^{2S+L_5}
 \le b_4^{m-3}(s_5/b_4)k^{r+p-4}.
\]
Now put $r=1+i$, $p=2+i+j$, $m=3+i+j+\ell$, with
$i,j,\ell\ge0$, and sum the three geometric series.

At $z=633/1000$ take
\[
 a_5=\frac{876716}{10^6},\quad b_4=\frac{982645}{10^6},\quad
 \tau=\frac{928171}{10^6}.
\]
Every norm comparison is an exact rational inequality, since
\[
 \operatorname{tr}D_t(z)^4=
 \sum_{i_1,\ldots,i_4=1}^t
 z^{i_1+i_2+i_3+i_4}
 \prod_{j=1}^4{\bf1}_{i_j+i_{j+1}\ge t},\qquad i_5=i_1.
\]
The two trace sums are strictly below $a_5^4,b_4^4$, respectively;
the endpoint inequalities hold, and
$1/2<\sqrt{a_5b_4}<\tau<1$.
The preceding bound equals
$z^{10}s_5/[\sqrt{a_5b_4}(1-\sqrt{a_5b_4})(1-a_5)(1-b_4)]$.
The function $1/[x(1-x)]$ increases for $1/2<x<1$, so replacing the
square root by $\tau$ gives the rational upper bound
\[
 \frac{z^{10}s_5}{\tau(1-\tau)(1-a_5)(1-b_4)}
 <112.194409824952<113.
\]
All summations can first be made over finite boxes; monotone convergence
then proves convergence as well as the stated bound.
\end{proof}

\subsection{Conditional self-pair contraction}

Write $Z_{m,r,p}(z)$ for the raw cap-four partition with its two
rightmost fours fixed at $r<p$, including their factor $z^8$.
Let $Z^{\rm self}_{m,r,p}(z)$ be the subpartition avoiding
\[
 F_i=\{x_i=1,\ x_{2i}=3\},\qquad p<i<m/2.
\]
Every genuine word belongs to this subpartition, by the noncarry Kunz
inequality at $(i,i)$. For $0<z\le1$ put
\begin{gather*}
 a=z,\quad b=z^2+z^3,\quad s=a+b,\qquad
 \epsilon=\frac{abz^3}{s(ab+s^2)},\quad
 \delta=\frac{\epsilon}{(1+z)^2},\\
 \beta=1-\epsilon,\quad \gamma=1-\delta,\quad T=\gamma/\beta,
 \quad d=p-r,\quad N=\max\{0,\lfloor(m-1)/2\rfloor-p\}.
\end{gather*}

\begin{lemma}\label{mass:contraction}
For every pivot pair and multiplicity,
\begin{equation}\label{mass:feedback}
 Z^{\rm self}_{m,r,p}(z)
 \le \beta^N T^{1+\lfloor d/3\rfloor}Z_{m,r,p}(z).
\end{equation}
For $m\le2p$ the stronger bound $Z^{\rm self}_{m,r,p}\le Z_{m,r,p}$
will be used.
\end{lemma}

\begin{proof}
Assume $m>2p$, and let
\[
 V=\{p+1,\ldots,\lfloor(m-1)/2\rfloor\},\qquad |V|=N.
\]
Give the raw words the probability measure proportional to their weights.
For a center $i\in V$, write
\[
 \ell_r(i)=m+r-i,\quad \ell_p(i)=m+p-i,\quad
 B_i=\{i,2i,\ell_r(i),\ell_p(i)\}.
\]
Every coordinate in $B_i$ exceeds $p$, so its alphabet is
$\{1,2,3\}$. Both leaves $\ell_r(i),\ell_p(i)$ exceed $m/2$.
The center has precisely these two retained neighbors, and has no loop.
All retained constraints touching the gadget are carry constraints of
threshold three. The only possible identification in $B_i$ is
$2i=\ell_r(i)$ or $2i=\ell_p(i)$.

\emph{Local failure families.}
First suppose the four vertices are distinct. Fix the coordinates outside
$B_i$ to an assignment admitting at least one completion. The family
\[
 x_i=1,\qquad x_{2i}=3,\qquad
 x_{\ell_r(i)},x_{\ell_p(i)}\in\{2,3\}
\]
has total gadget weight $ab^2z^3$. It satisfies all retained constraints
for every such exterior assignment: the center's constraints are
satisfied by its leaves, and every other gadget spin is at least two.
Dropping all restrictions except the two center--leaf edges bounds the
whole gadget partition by
\[
 s(ab^2+bs^2).
\]
Indeed center spin one contributes $ab^2$, center spin two or three
contributes at most $bs^2$, and the target contributes at most $s$.
The ratio of the displayed numerator to this denominator is $\epsilon$.
If the target coincides with one leaf, the corresponding numerator and
denominator are $abz^3$ and $ab^2+bs^2$. Their ratio is
\[
 \frac{az^3}{ab+s^2}>\epsilon.
\]
These bounds also allow any additional edge or loop among the gadget
vertices, since the failure families satisfy it and it only reduces
the denominator.

\emph{An ordering of the centers.}
For distinct centers draw $i\to j$ when
\begin{equation}\label{mass:dependency}
 2j\in\{\ell_r(i),\ell_p(i)\}.
\end{equation}
This edge means that $F_i$ should be processed before $F_j$ to retain
both choices at every reflection leaf. We construct a set of postponed
centers of size at most $1+\lfloor d/3\rfloor$, and an order of all
other centers in which every arrow goes to a later or postponed center.
Put
\[
 c=(2m+r+p)/6,\qquad I=[c-d/2,c+d/2].
\]
Process centers outside $I$ in decreasing order of $|i-c|$, with
arbitrary tie breaking. For an arrow $i\to j$, write
$j=(m+t-i)/2$ with $t\in\{r,p\}$. Since $m+t-3c=\pm d/2$,
\[
 |j-c|\le |i-c|/2+d/4<|i-c|\qquad(i\notin I).
\]
Thus an outside center never points to an earlier center. The weak
version of the same inequality shows that both affine maps preserve
$I$, so no center in $I$ points to an outside center.

The remaining centers form the consecutive integer set $V\cap I$,
with at most $d+1$ elements. If there are $d+1$ elements, postpone one
endpoint. The remaining set has at most $d$ elements and diameter
strictly less than $d$. For any fixed target vertex $j$, its two
possible incoming sources $m+r-2j$ and $m+p-2j$ differ by $d$;
therefore its indegree in the remaining graph is at most one.
Every nontrivial strongly connected component is consequently a simple
directed cycle: strong connectivity supplies an internal incoming edge
at every vertex, leaving exactly one internal incoming and outgoing
edge at each vertex. A two-cycle would give
\[
 i+2j=m+t_1,\quad j+2i=m+t_2,
 \qquad i-j=t_2-t_1\in\{-d,0,d\},
\]
which is impossible for distinct vertices of this set. Self-edges were
excluded in \eqref{mass:dependency}. Every directed cycle thus has at
least three vertices. Postpone one vertex from each cyclic component
and put the remaining acyclic graph in topological order, with $i$
before $j$ on every edge $i\to j$. The postponed set has size
\[
 f\le 1+\lfloor d/3\rfloor.
\]
All centers not postponed have now been ordered; process the postponed
ones last in any order.

\emph{Survival under the preceding avoidances.}
Consider a nonpostponed center $i$, and condition on avoidance of all
preceding $F_j$. On a fixed exterior assignment outside $B_i$ with
positive mass under that conditioning, we claim that the entire local
failure family above still survives. A preceding event disjoint from
$B_i$ is already avoided by the fixed exterior assignment. For one
meeting $B_i$, the possible intersections are exhaustive as follows.
An earlier center cannot be a reflection leaf, since centers are below
$m/2$ and leaves above it. If $j=2i$, its center has spin three, so
$F_j$ is false. If $2j=i$, its target has spin one, so $F_j$ is false.
Equal centers or equal doubled targets force $j=i$. The only remaining
possibility is a preceding target at a reflection leaf, precisely an
arrow $i\to j$, which the ordering excludes. In a three-vertex gadget,
a target at the shared leaf $2i$ would again force $j=i$.

It follows that the conditional numerator on this exterior slice is
at least the full weight already displayed. The conditional denominator
is at most the unrestricted gadget denominator displayed there.
The exterior weight cancels. Averaging over the exterior slices proves
\[
 \Pr(F_i\mid\text{all preceding avoidances})\ge\epsilon.
\]
This argument does not require preceding avoidances to be measurable
outside $B_i$; their survival was checked on each proposed completion.

For a postponed center use the single assignment
\[
 x_i=1,\qquad x_{2i}=3,\qquad
 x_{\ell_t(i)}=2\quad\text{when }\ell_t(i)\ne2i.
\]
It also makes false every preceding event meeting the gadget, including
an event whose target is a remaining leaf, now of spin two.
Its weight in the four-vertex case is $z^8$, giving conditional
probability at least
\[
 \frac{z^8}{s(ab^2+bs^2)}=\delta.
\]
In the three-vertex case its weight is $z^6$, and the bound is larger
because $s>z^2$. The same slice argument therefore applies in any order
to the postponed centers. Multiplying the conditional avoidance bounds
gives
\[
 Z^{\rm self}_{m,r,p}\le\beta^{N-f}\gamma^f Z_{m,r,p}
 \le\beta^N T^{1+\lfloor d/3\rfloor}Z_{m,r,p}.
\]
If any conditioning event has zero mass, the final left side is zero
and the assertion follows immediately. For $m\le2p$ there are no
eligible centers, and the raw bound suffices.
\end{proof}

\subsection{A uniform envelope for the absolute spectral expansion}

From now on set $z=z_4=1247/2000$, and define
\begin{gather*}
 D_{ij}=z^{(i+j)/2}{\bf1}_{i+j\ge4},\quad 1\le i,j\le4,\\
 C_{ij}=z^{(i+j)/2}{\bf1}_{i+j\ge3}{\bf1}_{i,j\le3},\quad
 J_3=\operatorname{diag}(1,1,1,0),\quad E=J_3DJ_3.
\end{gather*}
The nonzero carry eigenvalues are $\lambda>0>\mu$, where
$\lambda^2-b\lambda-ab=0$ and $\mu=b-\lambda$.
Let $Q=(q_+,q_-)$ have orthonormal eigenvectors as columns, with
$q_+$ nonnegative, and put $\rho=|\mu|/\lambda$. Thus
\[
 C^n=\lambda^nq_+q_+^T+\mu^nq_-q_-^T\qquad(n\ge1).
\]

We specify the spectral partitions used below. For a fixed reflection
graph with $m>2p$, retain each maximal strong run as an exact matrix product,
including its endpoint alphabet projections, and expand every maximal
carry run by the displayed identity. A mode assignment chooses $+$ or
$-$ for each carry run. After all strong blocks and boundaries have
been contracted, let $W_\omega$ denote its scalar weight, including
$z^8$ and any component having no carry edge. Define
\begin{equation}\label{mass:mode-definitions}
 P_{m,r,p}=W_{(+,\ldots,+)},\qquad
 A^{\rm abs}_{m,r,p}=\sum_\omega|W_\omega|,\qquad
 M_{m,r,p}=A^{\rm abs}_{m,r,p}-P_{m,r,p}.
\end{equation}
The positive mode has nonnegative endpoint and strong-block factors,
so $P_{m,r,p}\ge0$. Taking absolute values of the projected blocks and
endpoint factors before multiplying gives exactly $A^{\rm abs}_{m,r,p}$.
In particular
\begin{equation}\label{mass:spectral-majorization}
 0\le P_{m,r,p}\le A^{\rm abs}_{m,r,p},\qquad
 0\le Z_{m,r,p}\le P_{m,r,p}+M_{m,r,p}=A^{\rm abs}_{m,r,p}.
\end{equation}
Here the \emph{mixed} remainder $M_{m,r,p}$ includes every assignment
containing a negative mode, including the all-negative assignment.
This convention differs from the two-pure-mode decomposition used for
algebraic cancellation in Section~\ref{raw:topology}.

There is one carry component of zero ordinary-edge length that must
still receive two modes. It occurs at
$m=3p-2r>2p$, when deletion of port $p$ leaves the isolated vertex
$i=2p-r$ with a carry loop. Its free and loop boundary vectors are
\[
 h=(z^{1/2},z,z^{3/2},0)^T,\qquad
 h_+=(0,z,z^{3/2},0)^T.
\]
Both lie in $\operatorname{im}C$, so its scalar weight is split as
\begin{equation}\label{mass:zero-edge}
 h^Th_+=(h^Tq_+)(q_+^Th_+)+(h^Tq_-)(q_-^Th_+).
\end{equation}
This is the convention for $W_\omega$ at that corner; its absolute
mode sum is at most $\|h\|_2\|h_+\|_2$ by Cauchy--Schwarz.
An isolated vertex in the strong core is kept as its ordinary positive
scalar and is not given this artificial carry expansion.

The constants used below are
\begin{equation}\label{mass:constants}
 R=\frac{23389}{25000},\quad S=\frac{72789}{100000},\quad
 L=\frac{101778011}{10^8},\quad
 u=\begin{pmatrix}1\\7/25\end{pmatrix},\quad
 v=\begin{pmatrix}47/50\\9/50\end{pmatrix}.
\end{equation}

\begin{lemma}\label{mass:block}
We have $\|D\|_2<R$, $\|E\|_2<S$, $1<\lambda<L$ and
$\rho<39/100$. For every $k,n\ge1$, entrywise,
\begin{equation}\label{mass:actual-block}
 \left|Q^TD^kQ\operatorname{diag}(\lambda^n,\mu^n)\right|
 \le R^k\lambda^nuv^T.
\end{equation}
For a single restricted strong edge and every $n\ge1$,
\[
 \left|Q^TEQ\operatorname{diag}(\lambda^n,\mu^n)\right|
 \le S\lambda^nuv^T.
\]
\end{lemma}

\begin{proof}
Conjugating $D$ by $H=\operatorname{diag}(z^{j/2})_{j=1}^4$ gives
$M:=H^{-1}DH=(z^j{\bf1}_{i+j\ge4})_{ij}$. The positive test vector
$(1000000,1710962,2377416,2377416)^T$ satisfies $Mx<Rx$
componentwise. For the corresponding three-dimensional restriction $M_3$,
$(1000000,2146345,3002947)^T$ satisfies $M_3x<Sx$. Perron--Frobenius
and symmetry prove the norm bounds. The quadratic equation gives the
stated $\lambda,\rho$ inequalities by rational squaring.

Write $\alpha_k=q_+^TD^kq_+$, $\omega_k=q_+^TD^kq_-$ and
$\nu_k=q_-^TD^kq_-$. The following rational strict upper bounds on
absolute values suffice:
\[
\begin{array}{c|ccc}
 k&|\alpha_k|&|\omega_k|&|\nu_k|\\\hline
 1&68/100&18/100&5/100\\
 2&69/100&3/100&11/100\\
 4&\text{not needed}&\text{not needed}&2/100
\end{array}
\]
They can be checked without numerical eigenvectors. Set
\[
\begin{aligned}
 y_+&=(b/\lambda,1,1,0)^T,& y_-&=(\lambda/a,-1,-1,0)^T,\\
 \kappa_+&=a(b/\lambda)^2+b,& \kappa_-&=\lambda^2/a+b.
\end{aligned}
\]
Then $q_\pm=Hy_\pm/\sqrt{\kappa_\pm}$. Substitute the positive root in
these finite expressions, enclose its square root by rational squares,
and compare; the denominator signs are positive. This also gives the
same first-power moments for $E$, because $J_3q_\pm=q_\pm$.

For $k\ge2$, Cauchy--Schwarz gives
$|\alpha_k|/R^k\le(69/100)/R^2$. For $k\ge3$ it gives
\[
 \frac{|\omega_k|}{R^k}\le
 \frac{\sqrt{(69/100)(2/100)}}{R^3},\qquad
 \frac{|\nu_3|}{R^3}\le
 \frac{\sqrt{(11/100)(2/100)}}{R^3},
\]
and for $k\ge4$, $|\nu_k|/R^k\le(2/100)/R^4$.
Indeed apply $D$ to $q_+$ and $D^2$ to $q_-$ for the off-diagonal
estimate, and $D^2$ to both copies of $q_-$ for the last estimate.
For $k=1,2$ use the table directly. The second column in
\eqref{mass:actual-block}, after division by $\lambda^n$, gains
$\rho^n\le39/100$. Comparison with the four entries of $uv^T$ proves
the assertion; square the positive off-diagonal bounds if desired.
The same comparisons with denominator $S$ prove the restricted case.
This extends the finite moment checks to every power.
\end{proof}

The hypothesis $n\ge1$ is essential: the assertion is about a strong
run followed by a positive-length carry run. For example,
$\nu_2/R^2>3/25>(uv^T)_{22}$, so a bare strong-run assertion would
be false. The zero-edge component in \eqref{mass:zero-edge} is handled by its
boundary estimate; \eqref{mass:actual-block} is never used with $n=0$.

Here are the endpoint and component consequences of the lemma. A vertex
with index greater than $r$ has at most one strong incidence, so every
internal vertex of a strong run has the full alphabet $\{1,2,3,4\}$.
Alphabet restrictions on such a run therefore occur only at its endpoints.
At a carry junction they do not change the projected block, since
$J_3Q=Q$. Thus a run between carry components has the projected matrix
$Q^TD^kQ$; if both endpoints of a strong edge exceed $r$, that edge is
an entire one-edge strong run and the sharper $E$ estimate applies.
At a free or loop endpoint, absorb the alphabet and unary restrictions
into the boundary vector. Strong prefixes and suffixes are then bounded
by ordinary operator norms, using $\|E\|_2<S$ for the restricted
one-edge case. Alphabet projections and $Q^T$ are contractions, while
taking coordinatewise absolute values preserves Euclidean norm. Moreover
\[
 \|u\|_2^2\|v\|_2^2=\frac{77173}{78125}<1,\qquad
 v^Tu=\frac{619}{625}<1.
\]
Every ordinary cycle has both a strong and a carry edge: two successive
reflections in a cycle of either single kind translate an integer by the
nonzero distance $p-r$, and cannot close. Thus cycles can be written as
products of blocks with $k,n\ge1$. Products introduce no extra path
factor, and cycle traces are at most one after normalization. The at
most two free endpoint incidences together contribute at most $s_4=\sum_{j=1}^4z^j$; loop endpoint
norms are at most one, since $z^2+z^3+z^4<1<s_4$.
Counting strong edges gives
\[
 e_D=\lfloor(r-1)/2\rfloor+\lfloor(p-1)/2\rfloor
       -{\bf1}_{p\ne2r},\qquad
 e_E=\max\{\lfloor(p-1)/2\rfloor-r,0\}.
\]
With $\eta=R/L$ and $\sigma=S/R$, define
\begin{equation}\label{mass:raw-envelope}
 \mathcal E_{m,r,p}=z^8s_4L^{m-3}\eta^{e_D}\sigma^{e_E}.
\end{equation}
The argument proves $Z_{m,r,p}\le\mathcal E_{m,r,p}$ for every $m>p$,
and $A^{\rm abs}_{m,r,p}\le\mathcal E_{m,r,p}$ for every $m>2p$.
For short multiplicities the ordinary graph estimate uses the same
block bounds whenever there are carry edges and ordinary operator
norms otherwise.
The number of carry edges is at most $m-3-e_D$:
the graph has $m-3$ vertices and is a union of paths and cycles. The
restricted-edge replacement supplies $(S/R)^{e_E}$, while $L>1$
permits the stated larger carry exponent.

\subsection{Finite bases and complete multiplicity tails}

Write $f_d=1+\lfloor d/3\rfloor$. Combining
\eqref{mass:feedback} with \eqref{mass:spectral-majorization} gives the
nonnegative decomposition that will be summed:
\begin{equation}\label{mass:master-sum}
 A(z)\le
 \sum_{r<p}\sum_{m=p+1}^{2p}Z_{m,r,p}
 +\sum_{r<p}\sum_{m>2p}T^{f_d}\beta^N
             (P_{m,r,p}+M_{m,r,p}).
\end{equation}
It is enough initially to read this inequality over finite ranges;
the bounds below establish convergence. We keep the exact factor
$\beta^N$ in this decomposition. For the finite bases and mixed tails
we may enlarge it using the rational normalization
\begin{equation}\label{mass:w}
 w=\frac{70651823329321825}{72057594037927936},\qquad
 w^2\ge\beta,\quad w<1,\quad J=(wL)^2<1.
\end{equation}

Multiplying every vertex weight, including the two pivots, by $w$
gives the normalized partitions
\[
 \widehat P_{m,r,p}=w^{m-1}P_{m,r,p},\quad
 \widehat A^{\rm abs}_{m,r,p}=w^{m-1}A^{\rm abs}_{m,r,p},\quad
 \widehat M_{m,r,p}=w^{m-1}M_{m,r,p}.
\]
For $m>2p$ put
\[
 K_{m,r,p}=T^{f_d}w^{-2p-((m-1)\bmod2)}.
\]
Since $m-1-2p-((m-1)\bmod2)=2N$, for $X=P,M$ we have
\begin{equation}\label{mass:normalized-majorization}
 T^{f_d}\beta^N X_{m,r,p}
 \le K_{m,r,p}\widehat X_{m,r,p}.
\end{equation}
A normalized block in \eqref{mass:actual-block} gains $w^{k+n}$
on both sides; the boundary vectors each gain $w^{1/2}$ and the
pivot factor gains $w^2$. In particular the normalized zero-edge
component is $w h^Th_+$, with both terms of
\eqref{mass:zero-edge} multiplied by $w$.

The carry-port topology of Section~\ref{raw:topology} identifies the
carry components under $m\mapsto m+2d$: each paired-port component gains
four ordinary edges and each fixed-port component gains two, including
the component from which $p$ was deleted. The strong blocks and
boundary factors remain the same. Every assignment contributing to
$M$ uses a negative mode on at least one such component. Its absolute
weight therefore gains a factor at most
$(w\lambda)^{2d}\rho^2$ after normalization. Summing these termwise
bounds proves
\begin{equation}\label{mass:mixed-progression}
 \widehat M_{m+2d,r,p}
 \le (w\lambda)^{2d}\rho^2\widehat M_{m,r,p}.
\end{equation}

For the positive mode, the rank-one projector $q_+q_+^T$ supplies the
same vector $q_+$ at every surviving port. Contracting these vectors
against the fixed strong core is independent of which ports were
paired in the carry graph. The carry-loop endpoints each contribute
$q_+^Th_+$, and the free endpoint left by deleting $p$ contributes
$h^Tq_+$. The number of carry loops depends only on parity, and the
total carry exponent increases by two under $m\mapsto m+2$.
The zero-edge split in \eqref{mass:zero-edge} has exactly these factors.
Consequently
\begin{equation}\label{mass:pure-progression}
 \widehat P_{m+2,r,p}=(w\lambda)^2\widehat P_{m,r,p}
 \qquad(m>2p).
\end{equation}
Thus the two bases $2p+1,2p+2$ suffice for the pure sum, whereas the
$2d$ bases $2p+s$, $1\le s\le2d$, suffice for the mixed bound.
The multiplier $K_{m,r,p}$ is constant on each progression being used.

Here are explicit definitions of the three finite pivot sums. For a
positive integer cutoff $P_0$, set
\begin{align}
 \mathsf S(P_0)
 &=\sum_{p=2}^{P_0}\sum_{r=1}^{p-1}\sum_{m=p+1}^{2p}Z_{m,r,p},
 \label{mass:short-box}\\
 \mathsf P(P_0)
 &=\sum_{p=2}^{P_0}\sum_{r=1}^{p-1}\sum_{s=1}^{2}
 \frac{K_{2p+s,r,p}\widehat P_{2p+s,r,p}}
      {1-(w\lambda)^2},\label{mass:pure-box}\\
 \mathsf M(P_0)
 &=\sum_{p=2}^{P_0}\sum_{r=1}^{p-1}\sum_{s=1}^{2d}
 \frac{K_{2p+s,r,p}\widehat M_{2p+s,r,p}}
      {1-(w\lambda)^{2d}\rho^2}.\label{mass:mixed-box}
\end{align}
All denominators are positive because $w\lambda<wL<1$.
Equations~\eqref{mass:normalized-majorization}--\eqref{mass:pure-progression}
show exactly how these finite sums bound the corresponding complete
multiplicity sums in \eqref{mass:master-sum}.

For evaluating them one may use the right eigenvectors $y_\pm$ from
the proof of Lemma~\ref{mass:block}, with left vectors
$z^j y_{\pm,j}/\kappa_\pm$. Their positive diagonal change of basis
telescopes through complete products and endpoints and commutes with
entrywise absolute values. It therefore yields the same pure and
absolute-mode scalars just defined.

It remains to sum over pivots, rather than merely checking finitely
many of them. Write
\[
 r=2a+u_0,\quad p=2b+v_0,\quad u_0,v_0\in\{1,2\},\quad
 h=\begin{cases}0,&u_0<v_0,\\1,&u_0\ge v_0.\end{cases}
\]
The exact domain is $a\ge0$, $b\ge a+h$. Also
$e_D\ge a+b-1$ and $e_E=\max\{b-2a-u_0,0\}$. Define
\[
 \mathcal H_{u_0v_0}(A,B)
 =\sum_{a\ge0}\sum_{b\ge a+h}A^aB^b
      \sigma^{\max\{b-2a-u_0,0\}}.
\]
Splitting at $b=2a+u_0$ gives
\begin{equation}\label{mass:parity-sum}
 \mathcal H_{u_0v_0}(A,B)=
 \frac{B^h}{(1-B)(1-AB)}
 -\frac{B^{u_0+1}}{(1-B)(1-AB^2)}
 +\frac{B^{u_0+1}\sigma}{(1-B\sigma)(1-AB^2)}.
\end{equation}
The first two terms are the sum of finite inner strips; they do not
require $B<1$. The convergence requirements are
\begin{equation}\label{mass:convergence}
 AB<1,\qquad AB^2<1,\qquad B\sigma<1.
\end{equation}
At $B=1$ the removable expression is read by its limit, although none
of the evaluations below has $B=1$. Let
$\mathcal H^{>P}_{u_0v_0}$ be this rational expression minus the finite
sum over $2b+v_0\le P$, $0\le a\le b-h$. This is exactly the omitted
pivot sum, not an asymptotic approximation.

The short multiplicities $p+1\le m\le2p$, summed in
\eqref{mass:raw-envelope}, contribute for $p>P$ at most
\begin{equation}\label{mass:short}
 \frac{z^8s_4L^{-2}}{\eta(L-1)}
 \sum_{u_0,v_0}\left\{
 L^{2v_0}\mathcal H^{>P}_{u_0v_0}(\eta,\eta L^4)
 -L^{v_0}\mathcal H^{>P}_{u_0v_0}(\eta,\eta L^2)\right\}.
\end{equation}
Indeed the multiplicity sum of $L^{m-3}$ is
$L^{-2}(L^{2p}-L^p)/(L-1)$, and $\eta^{e_D}\le\eta^{a+b-1}$.

For the mixed tail put
\[
 t=\frac{252061464497}{250000000000},\qquad t^3\ge T,
 \qquad c_0=(39/100)^2,
\]
and define the rational function, on its convergence domain,
\begin{equation}\label{mass:G}
 G_P(x)=\sum_{u_0,v_0}L^{2v_0}x^{v_0-u_0}
 \mathcal H^{>P}_{u_0v_0}
       (\eta/x^2,\eta L^4x^2).
\end{equation}
This is the nonnegative sum over omitted pivots of
$L^{2p}\eta^{a+b}\sigma^{e_E}x^d$.
Pair the bases $s=2j+1,2j+2$, $0\le j<d$. Their normalization leaves
$w^{s-1-((s-1)\bmod2)}$, so their envelope is
$(1+L)L^{2p-2}J^j$. Thus the mixed multiplicity factor is
\[
 H_d=\frac{1-J^d}{(1-J)(1-c_0J^d)}.
\]
The following elementary bound retains the numerator that would be
lost by replacing its geometric sum with $d$.

\begin{lemma}\label{mass:three-evaluation}
For $0<J,c_0<1$, $a_d\ge0$ and $G(x)=\sum_{d\ge1}a_dx^d<\infty$,
\[
 \sum_{d\ge1}a_dx^dH_d\le
 \frac{G(x)-G(xJ)+\dfrac{c_0}{1-c_0}
              \{G(xJ)-G(xJ^2)\}}{1-J}.
\]
\end{lemma}
\begin{proof}
For $X=J^d$, use
$1/(1-c_0X)=1+c_0X/(1-c_0X)\le1+c_0X/(1-c_0)$,
multiply by $(1-X)/(1-J)$ and sum. Both differences are sums of
nonnegative terms of convergent series.
\end{proof}

Since $T^{1+\lfloor d/3\rfloor}\le Tt^d$, the mixed omitted tail is
at most
\begin{equation}\label{mass:mixed}
 \frac{C_0}{1-J}\left[
 G_P(t)-G_P(tJ)+\frac{c_0}{1-c_0}\{G_P(tJ)-G_P(tJ^2)\}\right],
 \qquad C_0=\frac{z^8s_4(1+L)L^{-2}T}{\eta}.
\end{equation}
All parameters here are rational. For $x=t,tJ,tJ^2$ the conditions
\eqref{mass:convergence} hold: passing from $x$ to $xJ$ leaves $AB$
unchanged and decreases $AB^2$ and $B\sigma$. The initial inequalities
are checked with the exact $w$ of \eqref{mass:w}, rather than replacing
$w^2$ by the slightly smaller $\beta$.

For the pure omitted sector, return to the exact factor $\beta^N$
in \eqref{mass:master-sum}. Equations~\eqref{mass:spectral-majorization}
and \eqref{mass:raw-envelope} give, for every $m>2p$,
\[
 T^{f_d}\beta^N P_{m,r,p}
 \le T^{f_d}\beta^N A^{\rm abs}_{m,r,p}
 \le T^{f_d}\beta^N\mathcal E_{m,r,p}.
\]
This argument does not require $P_{m,r,p}\le Z_{m,r,p}$, which need
not hold for a signed spectral expansion. To bound the sum, add the
nonnegative terms $p<m\le2p$ and use
$\beta^{\lfloor(m-1)/2\rfloor-p}$ throughout. On the added range its
exponent is nonpositive; on the original range the factor equals $\beta^N$.
Put
$J_0=L^2\beta<1$,
\[
 A_0=\eta/t^2,\quad B_0=\eta J_0\beta^{-2}t^2,\qquad
 M_1=L^{-2}J_0+L^{-1},\quad M_2=(L^{-2}+L^{-1})J_0.
\]
An upper bound for this enlarged sum with $p>P$ is
\begin{equation}\label{mass:full-feedback}
 z^8s_4T\sum_{u_0,v_0}
 \frac{\beta^{-v_0}t^{v_0-u_0}M_{v_0}}{\eta(1-J_0)}
 \mathcal H^{>P}_{u_0v_0}(A_0,B_0).
\end{equation}
For $p=2b+v_0$ the exact multiplicity sum needed here is
\[
 \sum_{m\ge p+1} L^{m-3}\beta^{\lfloor(m-1)/2\rfloor-p}
 =\beta^{-v_0}(J_0\beta^{-2})^b
     \frac{M_{v_0}}{1-J_0}.
\]
Substitute the pivot bounds to obtain \eqref{mass:full-feedback}.
Again $J_0<1$ and all three conditions in \eqref{mass:convergence}
hold. The short-tail parameters in \eqref{mass:short} also satisfy
those three conditions. Thus none of these expressions subtracts
divergent infinite sums.

\subsection{The finite interval calculation}

For the finite pivot prefixes the preceding formulas require finitely
many exact graph evaluations. For each fixed graph, path dynamic
programming sums over the endpoint spin, and a cycle is evaluated by
fixing its initial spin and summing the closing edge. Loop restrictions
are applied at the corresponding vertex. For pure and absolute-mode
partitions the same multiplication is performed on the projected
blocks. An upper enclosure of $\widehat A^{\rm abs}_{m,r,p}$ minus
a lower enclosure of $\widehat P_{m,r,p}$ is a valid upper enclosure
of $\widehat M_{m,r,p}$.

All interval endpoints are integers in units of $2^{-56}$. Each
multiplication and signed division is rounded outwards. The rational
$w$, a quadratic enclosure for $\lambda$, and a rational enclosure
for $T$ are checked before evaluation. Integer endpoints are guarded
by $2^{120}$; products that exceed signed 128-bit intermediate capacity
are evaluated with arbitrary-precision integers. The resulting upper
enclosures of the sums
\eqref{mass:short-box}--\eqref{mass:mixed-box} are as follows.
Here $U=2^{56}$. The input enclosures are
\[
 \frac{73338785529844359}{U}\le\lambda\le
 \frac{73338785529844360}{U},
\]\[
 \frac{73854862990399022}{U}\le T\le
 \frac{73854862990399023}{U}.
\]
The stored output fields \texttt{pure\_hi\_units},
\texttt{short\_hi\_units}, and \texttt{mixed\_hi\_units} are integer
upper endpoints in units of $U^{-1}$; the table displays their values
after division by $U$.

\begin{table}[htbp]
\centering
\caption{Finite pivot sums at $z=1247/2000$ with complete multiplicity
bounds. The integer numerators are scaled by $U=2^{56}$.}
\label{tab:mass-boxes}
\begin{tabular}{p{0.40\linewidth}r p{0.36\linewidth}}
\hline
Domain and evaluated quantity & Graph cases & Certified upper bound\\
\hline
$\mathsf P(400)$, two pure parity bases
 &159600&$101235082714477743831/U$\\
$\mathsf S(100)$ and $\mathsf M(100)$, short and mixed bases
 &666600&$\mathsf S(100)$: $3738038979785516612/U$;
 $\mathsf M(100)$: $36617965435406954/U$\\
\hline
\end{tabular}
\end{table}

The counts follow directly from
$2\sum_{p=2}^P(p-1)=P(P-1)$ and
$\sum_{p=2}^P\sum_{r=1}^{p-1}(p+2(p-r))
=2P(P+1)(P-1)/3$. The displayed bounds already include the respective
geometric multiplicity denominators and feedback multipliers, not just
the base-graph partitions. The finite graph evaluations establish the
displayed interval bounds;
the progression identities and inequalities above supply their extension
to every remaining multiplicity for these pivots.

Combining the boxes with \eqref{mass:short}, \eqref{mass:mixed} and
\eqref{mass:full-feedback} gives the following convenient outward
ceilings. Each finite decimal in this table denotes the indicated
rational number.

\begin{table}[htbp]
\centering
\caption{Rational upper bounds for the six terms covering the cap-four
mass. Every displayed decimal is a terminating rational value.}
\label{tab:mass-assembly}
\begin{tabular}{lr}
\hline
Contribution & Strict upper bound\\\hline
Pure prefix, $p\le400$ & $1404.918996618068$\\
Short prefix, $p\le100$ & $51.875711778803$\\
Mixed prefix, $p\le100$ & $0.508176354267$\\
Short omitted tail, $p>100$ & $116.098136035634$\\
Mixed omitted tail, $p>100$ & $446.302688776606$\\
Full feedback tail, $p>400$ & $3.745874303725$\\\hline
\end{tabular}
\end{table}

The sum of the displayed ceilings is below $2024$. Direct addition
before rounding gives the slightly sharper certified value
\[
 A(1247/2000)<2023.449583867101<2024.
\]
Every genuine word has a unique pair of rightmost maxima and is
covered by the contracted population. In \eqref{mass:master-sum}, the
short part is covered by $\mathsf S(100)$ and \eqref{mass:short}, the
pure part by $\mathsf P(400)$ and \eqref{mass:full-feedback}, and the
mixed part by $\mathsf M(100)$ and \eqref{mass:mixed}. All six bounds
are finite. Monotone convergence applied to the nonnegative sums proves
convergence and completes the proof of \eqref{mass:values}.

\section{Compensation and the final genus comparison}
\label{mass:final}

Lemma~\ref{cmb:compensation} combines the depth-four residual leaves
with the nonleaf parents of the unique-five residual leaves into disjoint
parts of the population counted by $A$. It is a positive-real weighted
comparison, not a coefficientwise comparison of equal genera. Substituting
\eqref{mass:values} in \eqref{cmb:mass-extraction} gives, for every $g$,
\begin{equation}\label{mass:residual}
 |\mathcal E_g|<2024z_4^{-g}+113z_5^{-g},\qquad
 z_4=1247/2000,\quad z_5=633/1000.
\end{equation}

\begin{proposition}\label{mass:tail}
For every source genus $g\ge837$,
$|\mathcal E_g|<L(g)$, where $L(g)$ is the unused-target capacity
of Proposition~\ref{cmb:capacity}.
\end{proposition}

\begin{proof}
Put
\[
 \bar q=\frac{309017}{500000},\qquad s_* =\frac{223607}{100000},
 \qquad r_* =\frac{809}{500}.
\]
The rational inequalities $\bar q^2+\bar q>1$, $s_*^2>5$ give
$\bar q>\varphi^{-1}$ and $s_*>\sqrt5$. Binet's formula therefore implies
\[
 \frac{|\mathcal E_g|}{\Fib_g}
 < \Xi(g):=\frac{s_*\{2024(\bar q/z_4)^g+113(\bar q/z_5)^g\}}
 {1-\bar q^{2g}}.
\]
The denominator uses the smaller $1-\varphi^{-2g}$ for both parities;
replacing $\varphi^{-1}$ by $\bar q$ only enlarges the upper bound.
Proposition~\ref{cmb:capacity} gives, throughout this range,
\[
 \frac{L(g)}{\Fib_g}\ge
 C(g):=1+\frac{927}{500}
       -\frac{50r_*^2}{((63/100)r_*)^g}.
\]
Exact rational evaluation at the joining genus gives
\[
 \Xi(837)<2.850575297189
 <2.853985748250<C(837).
\]
Since $\bar q<z_4,z_5$ and $(63/100)r_*>1$, the numerator terms of
$\Xi(g)$ strictly decrease, its denominator increases, and $C(g)$
increases. The comparison follows for every $g\ge837$. Failure of
this particular scalar comparison at $836$ has no implication for
monotonicity at that genus; it is covered by the finite comparison.
\end{proof}


\section{The finite comparison}\label{sec:computation}

We describe the exact finite calculation used with the preceding estimates.
All coefficients in this section are integers. The two pivots are the
rightmost occurrences of the maximum coordinate; this choice prevents
multiple counting when the pivot polynomials are summed.

Let $C_g(300)$ be the coefficient of $z^g$ in the sum of the raw
depth-four polynomials with rightmost pivot $p\le300$, over all
multiplicities and left pivots. Set
\begin{equation}\label{cmp:U}
 U_g=C_g(300)+\mathbf1_{\{g\ge604\}}
       \left\lfloor\frac{221715}{10^6}\Fib_g\right\rfloor.
\end{equation}
The omitted-pivot bound in Theorem~\ref{raw:tail} makes $U_g$ an upper
bound for the complete raw depth-four count. The support inequality for cap $a$,
\begin{equation}\label{cmp:support}
 2g\ge(a-1)m+p+1,
\end{equation}
also shows that $p>300$ cannot contribute below genus $604$ when $a=4$.
The floor is legitimate because the omitted coefficient is an integer.
Notice that this term first enters $U_g$ at $g=604$ and first enters the
shifted parent term $U_{g-1}$ at $g=605$.

For depth five, write $C^{(5)}_g$ for the coefficient of the complete raw
two-pivot population and define
\begin{equation}\label{cmp:V}
 V_g=\begin{cases}
 C^{(5)}_g,&0\le g\le300,\\
 \left\lfloor113(1000/633)^g\right\rfloor,&g\ge301.
 \end{cases}
\end{equation}
The first line is an exact coefficient calculation. The second line bounds
the genuine multiple-five population by the weighted mass estimate of the
preceding section; it is not a weighted-mass claim for the larger raw
population. Coefficients with negative indices are zero.

Let $N_g$ denote the no-one coefficients defined in the construction of the
partial injection. The residual bound is
\begin{equation}\label{cmp:residual}
 E^+_g=U_g-N_g-N_{g-3}-N_{g-4}
       +U_{g-1}-N_{g-1}+V_g.
\end{equation}
The shifted parent term has no sole-final-one subtraction: such a parent is
permitted by the unique-five map.

\begin{lemma}[Exact finite comparison]\label{cmp:finite}
For every $1\le g\le836$, the integers defined above satisfy
$0\le E^+_g\le L(g)$.
\end{lemma}

\begin{proof}[Verification and coverage]
For a truncation degree $H$, inequality~\eqref{cmp:support} bounds every
contributing multiplicity. The retained depth-four calculation has
$H=897$, $p\le300$, and
\[
 3\le m\le597,\qquad
 2\le p\le\min(300,m-1),\qquad 1\le r<p.
\]
These ranges comprise $17\,775\,550$ graph triples in $595$ disjoint
multiplicity classes. Each component polynomial is computed by the path or
cycle recurrence specified above, and the resulting products are summed and
truncated to give $C_g(300)$.

For depth five with $H=300$, the support inequality gives $m\le149$.
The sufficient complete run over $3\le m\le149$ and all pivot pairs visits
\[
 \sum_{m=3}^{149}\binom{m-1}{2}=\binom{149}{3}=540\,274
\]
graph triples. The resulting exact array contains all $301$ coefficients
through degree $300$. Thus no cap-five omitted-pivot estimate is used in the
exact range of~\eqref{cmp:V}.

The rational series defining $N_g$ and $L(g)$ are expanded by exact
coefficient recurrences. Substitution into
\eqref{cmp:U}--\eqref{cmp:residual} checks all $836$ inequalities. Put
$\Delta_g=L(g)-E^+_g$. The smallest integer slack is $1$, at $g=1,2$.
Among $1\le g\le836$, the smallest Fibonacci-normalized slack is
\[
 \min_g\frac{\Delta_g}{\Fib_g}
 =\frac{15028705}{33489287}>0.448761,
\]
attained at $g=42$. Table~\ref{tab:finite-transitions} records the switches
used by the proof. The displayed decimal entries are downward truncations of
exact rational values. The capacity-normalized slack also satisfies
\[
 \frac{\Delta_g}{L(g)}>0.179861\qquad(1\le g\le836),
\]
with its minimum attained at $g=605$.


\begin{table}[ht]
\centering\small
\caption{Transition checks for the finite comparison.}
\label{tab:finite-transitions}
\begin{tabular}{@{}rllr@{}}
\toprule
$g$ & depth-five term & depth-four tail status & $\Delta_g/\Fib_g$\\
\midrule
300 & exact $C^{(5)}_{300}$ & absent & $>0.871878$\\
301 & analytic mass bound & absent & $>0.683638$\\
603 & analytic mass bound & absent & $>0.872080$\\
604 & analytic mass bound & enters $U_g$ & $>0.650368$\\
605 & analytic mass bound & enters $U_{g-1}$ & $>0.513344$\\
836 & analytic mass bound & both terms present & $>0.513474$\\
\bottomrule
\end{tabular}
\end{table}


For completeness, the polynomial calculation introduces no numerical
approximation. A polynomial truncated at degree $H$ is evaluated at the
digit base $b=2^{H+2}$. The coefficient at actual weight $h\le H$ is at most
the number of positive compositions of $h$, hence less than $b$. Components
are placed in a fixed coordinate order, and the unique pivot choice preserves
this interpretation on summation. Minimum-weight normalization changes only
the position of a digit; a state is pruned only when its least possible
completed weight exceeds $H$. A negative normalization shift is made only
after an exact divisibility check. Discarded higher digits therefore cannot
alter a retained coefficient.

The exact transition rows, input identities, and regeneration procedure are
recorded in the repository~\cite{LiuCompanion}. Proposition~\ref{cmb:raw-envelope}
and the preceding analytic bounds identify these integer comparisons
with the required bound on $|\mathcal E_g|$.
\end{proof}

At the other end of the finite range,
$\Delta_{836}/\Fib_{836}>0.513474$. The scalar analytic comparison of
Proposition~\ref{mass:tail} fails at $836$, but at $837$ its exact rational
margin satisfies
\[
 C(837)-\Xi(837)>0.0034104510,
 \qquad \frac{C(837)-\Xi(837)}{C(837)}>0.0011949783.
\]
Thus the finite and analytic arguments meet without an uncovered genus.

\begin{proof}[Proof of Theorem~\ref{thm:main}]
For $g=0$, the two counts are $n_0=n_1=1$.
For $1\le g\le836$, the residual bound and
Lemma~\ref{cmp:finite} give $|\mathcal E_g|\le L(g)$.
For $g\ge837$, Proposition~\ref{mass:tail} gives the same comparison.
Applying~\eqref{cmb:fundamental} yields the result.
\end{proof}

\section{Concluding remarks}

The stronger Fibonacci inequality remains open:
\[
 n_{g+1}\ge n_g+n_{g-1}\qquad(g\ge1).
\]
The present construction guarantees enough unused targets
to accommodate its exceptional sources, but does not provide the
additional $n_{g-1}$ targets required for this inequality. The finer
error terms studied by Zhu~\cite{Zhu} offer a complementary way to
formulate that obstruction.

It would also be useful to replace the comparison of cardinalities
for the exceptional sources by a direct operation on their Kunz words,
with an explicit inverse on its image. The theorem already guarantees
an injection between the two finite genus classes. The question is
whether the remaining matching can be explained by a structural rule,
rather than by estimates and finite verification.


\section*{AI statement}
This work was developed through continuous collaboration between the
author and AI models. GPT-6 Astra (OpenAI) and Fable 5.1 (Anthropic)
contributed to proof development and checking, exact computation,
literature research, and preparation and review of the manuscript.
The author takes full responsibility for all mathematical results
and the contents of this manuscript.

\begin{thebibliography}{99}
\bibitem{Bacher} R. Bacher, \emph{Generic numerical semigroups},
arXiv:2105.04200v1, 2021.
\bibitem{Bras} M. Bras-Amor\'os, Fibonacci-like behavior of the number
of numerical semigroups of a given genus, \emph{Semigroup Forum}
\textbf{76} (2008), 379--384.
\bibitem{BT} M. Bras-Amor\'os and V. Torra,
\emph{Multiparameter counting of numerical semigroups: recurrences and
leaf-discriminating trees}, arXiv:2607.23111v3, 2026.
\bibitem{CMR} M. Casas, J. A. Madrid and J. C. Rosales,
\emph{Internal numerical semigroups}, arXiv:2608.19984v1, 2026.
\bibitem{CK} S. Cyrusian and N. Kaplan, Ordinarization numbers of
numerical semigroups, \emph{Communications in Algebra}, published
online 11 February 2026, \href{https://doi.org/10.1080/00927872.2026.2615092}
{doi:10.1080/00927872.2026.2615092}.
\bibitem{EF} S. Eliahou and J. Fromentin, Gapsets and numerical
semigroups, \emph{Journal of Combinatorial Theory, Series A}
\textbf{169} (2020), 105129.
\bibitem{Kaplan} N. Kaplan, Counting numerical semigroups,
\emph{American Mathematical Monthly} \textbf{124} (2017), no.~9, 862--875.
\bibitem{KR} A. Ataei Kachouei and F. Rahmati,
\emph{Gapset Extensions, Theory and Computations},
arXiv:2408.02425v1, 2024.
\bibitem{Kumar} A. Kumar, \emph{Numerical semigroups of ordinarization
number three}, preprint, 29 July 2026.
\url{https://www.researchgate.net/publication/410922576_NUMERICAL_SEMIGROUPS_OF_ORDINARIZATION_NUMBER_THREE}.
\bibitem{Kunz} E. Kunz, \emph{\"Uber die Klassifikation numerischer
Halbgruppen}, Regensburger Mathematische Schriften \textbf{11},
Universit\"at Regensburg, 1987.
\bibitem{LiuCompanion} M. Liu, \emph{Genus monotonicity for numerical
semigroups: code and data}, version v1, 2026.
\url{https://github.com/michaeliu4/bras-amoros-monotonicity/tree/26d999739c7233f8131790efcbc367c0463caf09}.
\bibitem{Rosales} J. C. Rosales, P. A. Garc\'ia-S\'anchez,
J. I. Garc\'ia-Garc\'ia and M. B. Branco, Systems of inequalities
and numerical semigroups, \emph{Journal of the London Mathematical
Society} \textbf{65} (2002), 611--623.
\bibitem{Zhai} A. Zhai, Fibonacci-like growth of numerical semigroups
of a given genus, \emph{Semigroup Forum} \textbf{86} (2013), 634--662.
\bibitem{Zhu} D. G. Zhu, Sub-Fibonacci behavior in numerical semigroup
enumeration, \emph{Combinatorial Theory} \textbf{3} (2023), no.~2,
Paper No.~10.
\end{thebibliography}

\clearpage
\appendix
\section{Computational material}\label{app:computation}

The computational part has three distinct tasks: generating the raw
coefficient arrays, bounding the finite weighted graph sums, and checking
the resulting scalar inequalities. Table~\ref{tab:computations}
records their parameter ranges. The verification of stored outputs and
the regeneration of those outputs are separate procedures.


\begin{table}[ht]
\centering\small
\caption{Finite inputs to the proof. The retained depth-four arrays
cover a wider range than the theorem requires.}
\label{tab:computations}
\begin{tabular}{@{}p{.37\textwidth}p{.36\textwidth}r@{}}
\toprule
Calculation & Range & Graphs\\
\midrule
Raw depth-four coefficients & $p\le300$, $3\le m\le597$;
degrees $0$--$897$ & $17\,775\,550$\\
Raw depth-five coefficients & $3\le m\le149$; all pivots;
degrees $0$--$300$ & $540\,274$\\
Pure-mode weighted sum & $p\le400$ & $159\,600$\\
Short and mixed weighted sums & $p\le100$ & $666\,600$\\
\bottomrule
\end{tabular}
\end{table}


The depth-five row describes sufficient regeneration coverage.
The retained file was originally produced over $3\le m\le293$
($4\,149\,466$ parameter triples); the support inequality makes all
additional contributions zero through degree $300$. Regeneration over
$3\le m\le149$ gives the same $301$ coefficients.

The coefficient calculation uses arbitrary-precision integer arithmetic.
The weighted graph sums use rational enclosures with outward rounding,
and the final scalar comparisons use exact rational arithmetic.
The calculation of an enclosure, its propagation through a recurrence,
and the final strict comparison are all specified by the mathematical
bounds in the text.

\paragraph{Code and data availability.}
The source code and exact finite outputs are available in the public
repository, version v1~\cite{LiuCompanion}. Its verification guide maps
the computations in the text to their inputs and commands, with separate
instructions for checking stored results and regenerating coefficient
arrays and interval bounds. All Python checks use the standard library;
regeneration requires a C++17 compiler with 128-bit integer support and
the GMP library.


\end{document}
